\documentclass[letterpaper,11pt]{exam}

\usepackage[activeacute,spanish]{babel}
\usepackage[left=1.8cm,top=1cm,right=1.8cm, bottom=1cm,letterpaper, includeheadfoot]{geometry}
\usepackage{framed}
\usepackage{wasysym}
\usepackage[utf8]{inputenc}
\usepackage{algorithmic}
\usepackage{algorithm}
\usepackage{enumerate}
\usepackage{multicol}
\usepackage{amssymb, amsmath, amsthm}
\usepackage{subcaption}
\usepackage{graphicx,txfonts}
\usepackage{lmodern,url}
\usepackage{wrapfig}
\usepackage{upgreek}
\usepackage{hyperref}
\usepackage[dvipsnames]{xcolor}
\usepackage{mdframed}
\usepackage{epigraph}
\usepackage{cancel}
\usepackage{tikz}
\usepackage{float}
\def\checkmark{\tikz\fill[scale=0.4](0,.35) -- (.25,0) -- (1,.7) -- (.25,.15) -- cycle;}
\floatname{algorithm}{Algoritmo}
\makeatletter

\usepackage{fancyhdr}
\setlength{\headheight}{14pt}
\pagestyle{fancy}
\fancypagestyle{plain}{%
    \fancyhf{}
    \lhead{\footnotesize\itshape\bfseries\rightmark}
    \rhead{\footnotesize\itshape\bfseries\leftmark}
    }
\renewcommand{\footrulewidth}{0.5pt}
\setlength{\parindent}{1cm}
\newenvironment{chapquote}[2][2em]
  {\setlength{\@tempdima}{#1}%
   \def\chapquote@author{#2}%
   \parshape 1 \@tempdima \dimexpr\textwidth-2\@tempdima\relax%
   \itshape}
  {\par\normalfont\hfill--\ \chapquote@author\hspace*{\@tempdima}\par\bigskip}
\makeatother


\newcommand{\arcangle}{%
  \mathord{<\mspace{-9mu}\mathrel{)}\mspace{2mu}}%
}

% macros
\newcommand{\heart}{\ensuremath\heartsuit}
\newcommand{\grad}{\hspace{-2mm}$\phantom{a}^{\circ}$}
\newcommand{\Q}{\mathbb Q}
\newcommand{\R}{\mathbb R}
\newcommand{\N}{\mathbb N}
\newcommand{\Z}{\mathbb Z}
\newcommand{\C}{\mathbb C}
\newcommand{\U}{\mathcal U}
\newcommand{\I}{\mathbb I}
\newcommand{\ssi}{\Longleftrightarrow}
\newcommand{\To}{\Rightarrow}
\newcommand{\tq}{\mid }
\newcommand{\exclusivo}{\veebar }
\renewcommand{\vec}[2]{\left(\begin{array}{c}{#1}\\{#2}\end{array}\right)}
\newcommand{\texii}[2]{\begin{minipage}{0.5\textwidth} #1 \end{minipage}
                     \begin{minipage}{0.5\textwidth} #2 \end{minipage}}

\providecommand{\abs}[1]{\lvert#1 \rvert}
\newcommand{\sech}{\operatorname{sech}}

% Teoremas, Lemas, etc.
\theoremstyle{plain}
\newtheorem{teo}{Teorema}
\newtheorem{lem}{Lema}
\newtheorem{con}{Conjetura}
\newtheorem{prop}{Proposici\'on}
\newtheorem{cor}{Corolario}
\newtheorem{prob}{Problema Controlable}
\newtheorem{nota}{Notaci\'on}
\newtheorem{obs}{Observaci\'on}
\newtheorem{definition}{Definici\'on}
\newtheorem{propiedades}{Propiedad}
\newtheorem{ax}{Axioma}
\newtheorem{ej}{Ejemplo}
\newtheorem{formula}{Fórmula}
\setlength{\parindent}{0pt}

%%%%%%% inicio documento %%%%%%%
\begin{document}

%============Encabezado estandar============

\newpage
\pagestyle{fancy}
\fancyhead[L]{\textit{Facultad de Ciencias Físicas y Matemáticas}}
\fancyhead[R]{\textit{Universidad de Chile}}

\noindent
\textbf{MA799-1  Introducción a la Tesis}\\
\textbf{Profesor: } Michal Kowalczyk\\
\textbf{Autor: } Isidora Miranda\\

\begin{center}
{\bf \Large Soluciones del sistema NLS acoplado}\\
\end{center}

\section{Introducción}
Un solitón es una onda solitaria que se propaga sin deformarse en un medio no lineal. El primero en observarla fue John Scott Russel (1808-1882), ingeniero escocés, al estudiar la propagación de ondas acuáticas en canales poco profundos.

\hspace{2mm}

Existen distintos tipos de solitones, sea $\psi$ la solución de un sistema luego

\begin{itemize}
    \item Bright Soliton: $|\psi|^2\to 0 $ cuando $x\to \pm \infty$. En este caso nuestras soluciones Bright corresponden a funciones de $\sech$.

\begin{figure}[H]
    \centering
    \includegraphics[width=0.6\textwidth]{TESIS/img/bright.png}
    \caption{Representación Bright Soliton}
    \label{fig:regiones_sol1}
\end{figure}
    
    \item Black Soliton: Sea $\rho$ una variable, tal que  la solución cumple que $|\psi|^2\to \rho $ cuando $x\to \pm \infty$ y alcanca su minimo en cero.En este caso nuestras soluciones Bright corresponden a funciones de $\tanh$.
    
\begin{figure}[H]
    \centering
    \includegraphics[width=0.6\textwidth]{TESIS/img/black.png}
    \caption{Representación Black Soliton}
    \label{fig:regiones_sol1}
\end{figure}
    \item Dark Soliton: Sea $\rho$ una variable, tal que  la solución cumple que $|\psi|^2\to \rho $ cuando $x\to \pm \infty$ y su minimo es mayor a cero. En este caso nuestras soluciones Bright corresponden a funciones de $\sech$ con una parte imaginaria.
    
\begin{figure}[H]
    \centering
    \includegraphics[width=0.6\textwidth]{TESIS/img/dark.png}
    \caption{Representación Dark Soliton}
    \label{fig:regiones_sol1}
\end{figure}
\end{itemize}




\section*{Contexto}
\vspace{-5mm}
\noindent\rule{\columnwidth}{0.4pt}
Sea el siguiente sistema de ecuaciones de derivadas parciales, donde trabajaremos en el espacio $(x,t)\in \R \times [0, +\infty)$
\begin{align}
    i \partial_t \psi&=\partial_{xx}\psi + \psi (1-|\psi|^2-\alpha|\phi|^2) \label{eq:NLS1} \\
     i \partial_t \phi&=\partial_{xx}\phi + \phi(1-\alpha|\psi|^2-\beta|\phi|^2) \label{eq:NLS2}
\end{align}

La motivación de nuestro trabajo es encontrar soluciones para este sistema, donde las determinaremos usando distintos enfoques y diversos casos.

Asociamos a este sistema el funcional de energía
\begin{equation}\label{eq:energia}
    E(\psi,\phi) = \int_{\mathbb{R}} \frac{1}{2}|\psi_x|^2 + \frac{1}{2}|\phi_x|^2
    + \frac{1}{4}(1 - |\psi|^2 - |\phi|^2)^2
    + \frac{\alpha - 1}{2}|\psi|^2|\phi|^2
    + \frac{\beta - 1}{4}|\phi|^4 \, dx,
\end{equation}
el cual se descompone naturalmente como $E = E_{\mathrm{cin}} + E_{\mathrm{pot}}$, donde
\begin{align}
    E_{\mathrm{cin}} &= \int_{\mathbb{R}} \frac{1}{2}|\psi_x|^2 + \frac{1}{2}|\phi_x|^2 \, dx, \\
    E_{\mathrm{pot}} &= \int_{\mathbb{R}} \frac{1}{4}(1 - |\psi|^2 - |\phi|^2)^2
    + \frac{\alpha-1}{2}|\psi|^2|\phi|^2 + \frac{\beta-1}{4}|\phi|^4 \, dx.
\end{align}
A continuación mostramos que $E$ es una cantidad conservada del sistema, es decir, $\partial_t E = 0$ cuando
$\psi$ y $\phi$ satisfacen \eqref{eq:NLS1}--\eqref{eq:NLS2}.

\subsection*{Derivación de la energía cinética}

Multiplicamos la ecuación \eqref{eq:NLS1} por $\overline{\psi_t}$ y tomamos parte real:
\begin{equation}
    \mathrm{Re}(i|\psi_t|^2) = \mathrm{Re}(\psi_{xx}\overline{\psi_t})
    + \mathrm{Re}\!\left(\psi(1 - |\psi|^2 - \alpha|\phi|^2)\overline{\psi_t}\right).
\end{equation}
Dado que $i|\psi_t|^2$ es imaginario puro, su parte real es cero y obtenemos
\begin{equation}
    0 = \mathrm{Re}(\psi_{xx}\overline{\psi_t})
    + \mathrm{Re}\!\left(\psi(1 - |\psi|^2 - \alpha|\phi|^2)\overline{\psi_t}\right).
\end{equation}
Trabajamos el primer término usando la identidad de producto
\begin{equation}
    \mathrm{Re}(\psi_{xx}\overline{\psi_t})
    = \partial_x \mathrm{Re}(\psi_x\overline{\psi_t}) - \mathrm{Re}(\psi_x\overline{\psi_{xt}}).
\end{equation}
Para el segundo sumando usamos que para cualquier función compleja $f$ se tiene
$\mathrm{Re}(f\bar{f}_t) = \frac{1}{2}\partial_t|f|^2$,
de modo que con $f = \psi_x$:
\begin{equation}
    \mathrm{Re}(\psi_x\overline{\psi_{xt}}) = \frac{1}{2}\partial_t|\psi_x|^2.
\end{equation}
Por lo tanto
\begin{equation}
    \mathrm{Re}(\psi_{xx}\overline{\psi_t})
    = \partial_x\mathrm{Re}(\psi_x\overline{\psi_t}) - \frac{1}{2}\partial_t|\psi_x|^2.
\end{equation}
Al integrar en $x \in \mathbb{R}$, el primer término es una divergencia espacial que se anula bajo las
condiciones de contorno $|\psi| \to \mathrm{cte}$ cuando $|x| \to \infty$. Así
\begin{equation}
    \int_{\mathbb{R}} \mathrm{Re}(\psi_{xx}\overline{\psi_t})\,dx
    = -\frac{1}{2}\partial_t\int_{\mathbb{R}}|\psi_x|^2\,dx.
\end{equation}
El procedimiento es completamente análogo para $\phi$, y al sumar ambas contribuciones obtenemos
\begin{equation}
    \int_{\mathbb{R}}\!\left[\mathrm{Re}(\psi_{xx}\overline{\psi_t})
    + \mathrm{Re}(\phi_{xx}\overline{\phi_t})\right]dx
    = -\partial_t E_{\mathrm{cin}}.
\end{equation}
Esto muestra que el término $\frac{1}{2}|\psi_x|^2 + \frac{1}{2}|\phi_x|^2$ es precisamente la
\emph{densidad de energía cinética}: la cantidad cuya derivada temporal aparece de forma natural
al multiplicar las ecuaciones por $\overline{\psi_t}$ y $\overline{\phi_t}$.

\subsection*{Derivación de la energía potencial}

Denotamos $\rho = |\psi|^2$ y $\sigma = |\phi|^2$ para simplificar la notación.
Los términos potenciales que resultan de multiplicar las ecuaciones por los respectivos conjugados son
\begin{equation}
    \mathrm{Re}\!\left(\psi(1-\rho-\alpha\sigma)\overline{\psi_t}\right)
    + \mathrm{Re}\!\left(\phi(1-\alpha\rho-\beta\sigma)\overline{\phi_t}\right).
\end{equation}
Como los factores $(1-\rho-\alpha\sigma)$ y $(1-\alpha\rho-\beta\sigma)$ son reales, y usando
$\mathrm{Re}(\psi\overline{\psi_t}) = \frac{1}{2}\partial_t\rho$, esto es igual a
\begin{equation}
    \frac{1}{2}(1-\rho-\alpha\sigma)\,\partial_t\rho
    + \frac{1}{2}(1-\alpha\rho-\beta\sigma)\,\partial_t\sigma.
\end{equation}
Expandimos cada factor y reagrupamos los términos según su naturaleza.

\medskip
\noindent\textbf{Términos lineales en $\rho$ y $\sigma$:}
\begin{equation}
    \frac{1}{2}\partial_t\rho + \frac{1}{2}\partial_t\sigma
    = \partial_t\!\left(\frac{\rho + \sigma}{2}\right).
\end{equation}

\noindent\textbf{Términos cuadráticos propios:}
\begin{equation}
    -\frac{1}{2}\rho\,\partial_t\rho - \frac{\beta}{2}\sigma\,\partial_t\sigma
    = -\frac{1}{4}\partial_t\rho^2 - \frac{\beta}{4}\partial_t\sigma^2
    = -\partial_t\!\left(\frac{|\psi|^4}{4} + \frac{\beta}{4}|\phi|^4\right).
\end{equation}

\noindent\textbf{Términos cruzados:}
\begin{equation}
    -\frac{\alpha}{2}\sigma\,\partial_t\rho - \frac{\alpha}{2}\rho\,\partial_t\sigma
    = -\frac{\alpha}{2}\partial_t(\rho\sigma)
    = -\frac{\alpha}{2}\partial_t(|\psi|^2|\phi|^2).
\end{equation}

Sumando los tres grupos, los términos potenciales se escriben como
\begin{equation}
    \partial_t\!\left(\frac{\rho+\sigma}{2} - \frac{|\psi|^4}{4}
    - \frac{\beta}{4}|\phi|^4 - \frac{\alpha}{2}|\psi|^2|\phi|^2\right).
\end{equation}
Completamos el cuadrado usando la identidad
\begin{equation}
    \frac{1}{4}(1-\rho-\sigma)^2
    = \frac{1}{4} - \frac{\rho+\sigma}{2} + \frac{\rho^2}{4} + \frac{\sigma^2}{4}
    + \frac{\rho\sigma}{2},
\end{equation}
de donde se obtiene
\begin{equation}
    \frac{\rho+\sigma}{2} - \frac{\rho^2}{4} - \frac{\sigma^2}{4} - \frac{\rho\sigma}{2}
    = \frac{1}{4} - \frac{1}{4}(1-\rho-\sigma)^2.
\end{equation}
Sustituyendo, y absorbiendo los términos con $(\alpha-1)$ y $(\beta-1)$:
\begin{align}
    &\frac{\rho+\sigma}{2} - \frac{|\psi|^4}{4} - \frac{\beta}{4}|\phi|^4
    - \frac{\alpha}{2}|\psi|^2|\phi|^2 \nonumber\\
    &\quad= \frac{1}{4} - \frac{1}{4}(1-\rho-\sigma)^2
    - \frac{\beta-1}{4}|\phi|^4 - \frac{\alpha-1}{2}|\psi|^2|\phi|^2.
\end{align}
La constante $\frac{1}{4}$ desaparece al tomar $\partial_t$, y al integrar en $x$ obtenemos
\begin{equation}
    \int_{\mathbb{R}}\!\left[\mathrm{Re}(\ldots)\right]dx = -\partial_t E_{\mathrm{pot}},
\end{equation}
donde
\begin{equation}
    E_{\mathrm{pot}} = \int_{\mathbb{R}} \frac{1}{4}(1-|\psi|^2-|\phi|^2)^2
    + \frac{\alpha-1}{2}|\psi|^2|\phi|^2 + \frac{\beta-1}{4}|\phi|^4\,dx.
\end{equation}

\subsection*{Conservación de la energía}

Juntando las contribuciones cinética y potencial, la identidad completa integrada es
\begin{equation}
    0 = -\partial_t E_{\mathrm{cin}} - \partial_t E_{\mathrm{pot}} = -\partial_t E(\psi,\phi),
\end{equation}
lo que prueba que $\partial_t E = 0$, es decir, $E$ es una cantidad conservada del sistema
\eqref{eq:NLS1}--\eqref{eq:NLS2}. El funcional de energía queda expresado como
\begin{equation}
    E(\psi,\phi) = \int_{\mathbb{R}} \frac{1}{2}|\psi_x|^2 + \frac{1}{2}|\phi_x|^2
    + \frac{1}{4}(1 - |\psi|^2 - |\phi|^2)^2
    + \frac{\alpha-1}{2}|\psi|^2|\phi|^2
    + \frac{\beta-1}{4}|\phi|^4 \, dx.
\end{equation}

\hspace{2mm}
\subsection*{Estructura Hamiltoniana}

El sistema \eqref{eq:NLS1}--\eqref{eq:NLS2} posee una estructura Hamiltoniana natural.
Decimos que un sistema de ecuaciones en derivadas parciales tiene estructura Hamiltoniana
si existe un funcional $E[\psi, \phi]$ tal que las ecuaciones de movimiento pueden escribirse
en la forma
\begin{equation}\label{eq:hamilton}
    i\partial_t\psi = -2\,\frac{\delta E}{\delta\bar{\psi}}, \qquad
    i\partial_t\phi = -2\,\frac{\delta E}{\delta\bar{\phi}},
\end{equation}
donde $\frac{\delta E}{\delta\bar{\psi}}$ denota la derivada funcional de $E$ respecto
a $\bar{\psi}$. Para un funcional de la forma $E = \int \mathcal{H}(\phi,\bar{\psi},\phi_x,\bar{\psi}_x)\,dx$,
esta derivada se calcula mediante la fórmula de Euler-Lagrange
\begin{equation}\label{eq:EL}
    \frac{\delta E}{\delta\bar{\psi}}
    = \frac{\partial\mathcal{H}}{\partial\bar{\psi}}
    - \partial_x\!\left(\frac{\partial\mathcal{H}}{\partial\bar{\psi}_x}\right),
\end{equation}
y análogamente para $\bar{\phi}$. El factor $-2$ en \eqref{eq:hamilton} es una convención
que emerge naturalmente de la forma del funcional, como veremos a continuación.

La motivación de esta estructura es que $E$ es una cantidad conservada del sistema:
si \eqref{eq:hamilton} se cumple, entonces $\frac{d}{dt}E = 0$ automáticamente,
como mostraremos al final de esta sección.

\subsection*{Construcción de $E$ a partir de las ecuaciones}

Para encontrar $E$, procedemos en sentido inverso: buscamos un funcional tal que
\eqref{eq:hamilton} reproduzca exactamente las ecuaciones \eqref{eq:NLS1}--\eqref{eq:NLS2}.
Esto equivale a requerir que
\begin{equation}
    -2\,\frac{\delta E}{\delta\bar{\psi}}
    = \partial_{xx}\psi + \psi(1 - |\psi|^2 - \alpha|\phi|^2).
\end{equation}
Identificamos dos tipos de contribuciones en el lado derecho: los términos de segundo orden
espacial $\partial_{xx}\psi$ y $\partial_{xx}\phi$, que llamaremos \emph{cinéticos}, y los
términos no lineales, que llamaremos \emph{potenciales}. Construimos $E = E_{\mathrm{cin}} + E_{\mathrm{pot}}$
tratando cada tipo por separado.

\medskip
\noindent\textbf{Energía cinética.}
Buscamos $E_{\mathrm{cin}}$ tal que $-2\frac{\delta E_{\mathrm{cin}}}{\delta\bar{\psi}} = \partial_{xx}\psi$.
Proponemos
\begin{equation}
    E_{\mathrm{cin}} = \int_{\mathbb{R}} \frac{1}{2}|\psi_x|^2 + \frac{1}{2}|\phi_x|^2\,dx.
\end{equation}
Verificamos usando \eqref{eq:EL}. Para el término $\frac{1}{2}|\psi_x|^2 = \frac{1}{2}\psi_x\bar{\psi}_x$:
\begin{equation}
    \frac{\partial}{\partial\bar{\psi}}\!\left(\frac{1}{2}|\psi_x|^2\right) = 0, \qquad
    \frac{\partial}{\partial\bar{\psi}_x}\!\left(\frac{1}{2}|\psi_x|^2\right) = \frac{1}{2}\psi_x,
\end{equation}
donde en la primera usamos que $|\psi_x|^2$ no depende de $\bar{\psi}$ directamente,
y en la segunda tratamos $\psi_x$ como constante según el truco de Wirtinger.
Aplicando \eqref{eq:EL}:
\begin{equation}
    \frac{\delta E_{\mathrm{cin}}}{\delta\bar{\psi}}
    = 0 - \partial_x\!\left(\frac{1}{2}\psi_x\right) = -\frac{1}{2}\psi_{xx}.
\end{equation}
Entonces:
\begin{equation}
    -2\,\frac{\delta E_{\mathrm{cin}}}{\delta\bar{\psi}} = -2\cdot\left(-\frac{1}{2}\psi_{xx}\right) = \psi_{xx}. \checkmark
\end{equation}
El factor $\frac{1}{2}$ en $E_{\mathrm{cin}}$ y el factor $-2$ en la convención
\eqref{eq:hamilton} no son independientes: es precisamente esta combinación la que
reproduce el término $\partial_{xx}\psi$ sin coeficientes adicionales. El procedimiento
es completamente análogo para $\phi$.

\medskip

\noindent\textbf{Energía potencial.}
Buscamos $E_{\mathrm{pot}}$ tal que
$-2\frac{\delta E_{\mathrm{pot}}}{\delta\bar{\psi}} = \psi(1 - |\psi|^2 - \alpha|\phi|^2)$.
Esto equivale a requerir
\begin{equation}
    \frac{\delta E_{\mathrm{pot}}}{\delta\bar{\psi}}
    = \frac{\psi}{2}\left(|\psi|^2 + \alpha|\phi|^2 - 1\right).
\end{equation}
Para encontrar $E_{\mathrm{pot}}$, observamos que el lado derecho contiene tres términos.
Calculamos qué funcional genera cada uno mediante \eqref{eq:EL}, usando que
$\frac{\partial|\psi|^2}{\partial\bar{\psi}} = \psi$ y $\frac{\partial|\psi|^4}{\partial\bar{\psi}} = 2|\psi|^2\psi$:
\begin{align}
    \frac{\delta}{\delta\bar{\psi}}\int \frac{|\psi|^4}{4}\,dx &= \frac{|\psi|^2\psi}{2},\\
    \frac{\delta}{\delta\bar{\psi}}\int \frac{\alpha}{2}|\psi|^2|\phi|^2\,dx &= \frac{\alpha}{2}|\phi|^2\psi,\\
    \frac{\delta}{\delta\bar{\psi}}\int \left(-\frac{|\psi|^2}{2}\right)dx &= -\frac{\psi}{2}.
\end{align}
Sumando estas tres contribuciones obtenemos exactamente
$\frac{\psi}{2}(|\psi|^2 + \alpha|\phi|^2 - 1)$, de modo que
\begin{equation}
    E_{\mathrm{pot}} = \int_{\mathbb{R}}
    \frac{|\psi|^4}{4} + \frac{\alpha}{2}|\psi|^2|\phi|^2 - \frac{|\psi|^2}{2}
    + \text{(términos en }\phi\text{)}\,dx.
\end{equation}
Los términos en $\phi$ se obtienen análogamente de la ecuación \eqref{eq:NLS2}, y resultan
ser $\frac{\beta}{4}|\phi|^4 - \frac{|\phi|^2}{2}$. Juntando todo:
\begin{equation}
    E_{\mathrm{pot}} = \int_{\mathbb{R}}
    - \frac{|\psi|^2}{2} - \frac{|\phi|^2}{2}
    + \frac{|\psi|^4}{4} + \frac{\alpha}{2}|\psi|^2|\phi|^2
    + \frac{\beta}{4}|\phi|^4\,dx.
\end{equation}
Esta expresión puede escribirse de forma más compacta. Usando la identidad
\begin{equation}\label{eq:expansion}
    \frac{1}{4}(1 - |\psi|^2 - |\phi|^2)^2
    = \frac{1}{4} - \frac{|\psi|^2}{2} - \frac{|\phi|^2}{2}
    + \frac{|\psi|^4}{4} + \frac{|\psi|^2|\phi|^2}{2} + \frac{|\phi|^4}{4},
\end{equation}
se verifica que
\begin{align}
    &- \frac{|\psi|^2}{2} - \frac{|\phi|^2}{2}
    + \frac{|\psi|^4}{4} + \frac{\alpha}{2}|\psi|^2|\phi|^2 + \frac{\beta}{4}|\phi|^4
    \nonumber\\
    &\quad= \frac{1}{4}(1-|\psi|^2-|\phi|^2)^2
    + \frac{\alpha-1}{2}|\psi|^2|\phi|^2
    + \frac{\beta-1}{4}|\phi|^4 - \frac{1}{4}.
\end{align}
La constante $-\frac{1}{4}$ no contribuye al funcional de energía (es independiente de
$\psi$ y $\phi$, por lo que su derivada funcional es cero), y la absorbemos escribiendo
la energía potencial en su forma compacta final:
\begin{equation}
    E_{\mathrm{pot}} = \int_{\mathbb{R}}
    \frac{1}{4}(1-|\psi|^2-|\phi|^2)^2
    + \frac{\alpha-1}{2}|\psi|^2|\phi|^2
    + \frac{\beta-1}{4}|\phi|^4\,dx.
\end{equation}
La forma compacta no es arbitraria: es exactamente el cuadrado que agrupa de forma
natural todos los términos no lineales que emergen de las ecuaciones.

\subsection*{Funcional de energía}

Sumando las contribuciones cinética y potencial, el funcional de energía asociado
al sistema \eqref{eq:NLS1}--\eqref{eq:NLS2} es
\begin{equation}\label{eq:Efinal}
    E(\psi,\phi) = \int_{\mathbb{R}} \frac{1}{2}|\psi_x|^2 + \frac{1}{2}|\phi_x|^2
    + \frac{1}{4}(1-|\psi|^2-|\phi|^2)^2
    + \frac{\alpha-1}{2}|\psi|^2|\phi|^2
    + \frac{\beta-1}{4}|\phi|^4\,dx.
\end{equation}

\subsection*{Conservación de la energía}

Mostramos que $\frac{d}{dt}E = 0$ cuando $\psi$ y $\phi$ satisfacen
\eqref{eq:NLS1}--\eqref{eq:NLS2}. Calculamos:
\begin{equation}
    \frac{d}{dt}E
    = \int_{\mathbb{R}} \frac{\delta E}{\delta\psi}\,\partial_t\psi
    + \frac{\delta E}{\delta\bar{\psi}}\,\partial_t\bar{\psi}
    + \frac{\delta E}{\delta\phi}\,\partial_t\phi
    + \frac{\delta E}{\delta\bar{\phi}}\,\partial_t\bar{\phi}\,dx.
\end{equation}
Usando que $\frac{\delta E}{\delta\psi} = \overline{\frac{\delta E}{\delta\bar{\psi}}}$
y agrupando en pares conjugados:
\begin{equation}
    \frac{d}{dt}E
    = 2\,\mathrm{Re}\!\int_{\mathbb{R}}
    \frac{\delta E}{\delta\bar{\psi}}\,\partial_t\bar{\psi}
    + \frac{\delta E}{\delta\bar{\phi}}\,\partial_t\bar{\phi}\,dx.
\end{equation}
Tomando complejo conjugado de la convención \eqref{eq:hamilton},
$i\partial_t\psi = -2\frac{\delta E}{\delta\bar\psi}$, obtenemos
$-i\partial_t\bar\psi = -2\frac{\delta E}{\delta\psi}$, es decir
$\partial_t\bar\psi = -2i\,\frac{\delta E}{\delta\psi}
= -2i\,\overline{\frac{\delta E}{\delta\bar\psi}}$.
Sustituyendo:
\begin{equation}
    \frac{d}{dt}E
    = 2\,\mathrm{Re}\!\int_{\mathbb{R}}
    \frac{\delta E}{\delta\bar{\psi}}\cdot\left(-2i\,\overline{\frac{\delta E}{\delta\bar\psi}}\right)
    + \frac{\delta E}{\delta\bar{\phi}}\cdot\left(-2i\,\overline{\frac{\delta E}{\delta\bar\phi}}\right)dx.
\end{equation}
Como $z\cdot(-2i\bar{z}) = -2i|z|^2$ es imaginario puro para cualquier $z$ complejo,
su parte real es cero. Por lo tanto:
\begin{equation}
    \frac{d}{dt}E = 0. \qquad \checkmark
\end{equation}
Esto confirma que $E$ definido en \eqref{eq:Efinal} es una cantidad conservada del
sistema \eqref{eq:NLS1}--\eqref{eq:NLS2}.

\newpage

\section{Solución 1}
En este caso usaremos ondas viajeras de la siguiente forma
\begin{align*}
    \psi(t,x)&=u(x-ct)\\
    \phi(t,x)&=v(x-ct)e^{i[(\lambda+\frac{c^2}{4})t-c\frac{x}{2}]}
\end{align*}

\textbf{Cálculo de derivadas ecuación \eqref{eq:NLS1}}

\begin{itemize}
    \item $i\partial_t \psi=-i c \dot{u}$
    \item $\partial_x \psi=\dot{u}$
    \item $\partial_{xx} \psi=\ddot{u}$
\end{itemize}

Luego obtenemos que la primera ecuación nos queda como
\begin{equation*}
    ic\dot{u}+\ddot{u}+u(1-|u|^2-\alpha|v|^2)=0.
\end{equation*}

\textbf{Cálculo de derivadas ecuación \eqref{eq:NLS2}}

\begin{itemize}
    \item $i\partial_t\phi=-ic\dot{v}e^{i[(\lambda+\frac{c^2}{4})t-c\frac{x}{2}]}-(\lambda+\frac{c^2}{4})ve^{i[(\lambda+\frac{c^2}{4})t-c\frac{x}{2}]}$
    \item $\partial_x\phi=\dot{v}e^{i[(\lambda+\frac{c^2}{4})t-c\frac{x}{2}]}+(\frac{-ic}{2})ve^{i[(\lambda+\frac{c^2}{4})t-c\frac{x}{2}]}$
    \item $\partial_{xx} \phi
    = \left(\ddot{v} - ic\dot{v} - \tfrac{c^2}{4}v \right)
    e^{i\left[(\lambda+\frac{c^2}{4})t-c\frac{x}{2}\right]}$
\end{itemize}

Donde obtenemos que
\begin{align*}
    i\partial_{t}\phi-\partial_{xx}\phi=-\lambda ve^{i[(\lambda+\frac{c^2}{4})t-c\frac{x}{2}]}-\ddot{v}e^{i[(\lambda+\frac{c^2}{4})t-c\frac{x}{2}]}
\end{align*}

Usando lo anterior, nuestro \textbf{sistema de ecuaciones} queda expresado de la siguiente forma
\begin{align*}
     ic\dot{u}+\ddot{u}&+u(1-|u|^2-\alpha|v|^2)=0 \\
     (\lambda v +\ddot{v})e^{i[(\lambda+\frac{c^2}{4})t-c\frac{x}{2}]}&+v e^{i[(\lambda+\frac{c^2}{4})t-c\frac{x}{2}]}(1-\alpha |u|^2-\beta|v|^2)=0
\end{align*}

Simplificamos, usando que $e^{i[(\lambda+\frac{c^2}{4})t-c\frac{x}{2}]}\neq0$ y finalmente obtenemos que
  \begin{align*}
     ic\dot{u}+\ddot{u}+u(1-|u|^2-\alpha|v|^2)&=0 \\
     (\lambda v +\ddot{v})+v (1-\alpha |u|^2-\beta|v|^2)&=0
\end{align*}

Reagrupamos términos y obtenemos
\begin{align}
     ic\dot{u}+\ddot{u}&+u(1-|u|^2-\alpha|v|^2)=0 \label{eq:u_sol1} \\
    \ddot{v}&+v (1+\lambda-\alpha |u|^2-\beta|v|^2)=0 \label{eq:v_sol1}
\end{align}

\subsection{Solución 1.1}\label{sec:sol11}
Proponemos los siguientes resultados
\begin{align*}
    u&=A_1 \tanh(bx)\\
    v&=A_2 \sech(bx)
\end{align*}

Calcularemos las derivadas de estas soluciones para obtener condiciones de existencia. Recordamos que $\sech'(x)=-\tanh(x)\cdot \sech(x)$.

\begin{itemize}
    \item $\dot{u}=A_1 b \sech^2(bx)$
    \item $\ddot{u}=A_1 b^2 \cdot 2 \sech(bx)(-\tanh(bx)\cdot \sech(bx)) = -2ub^2\cdot \sech^2(bx)$
\end{itemize}

Así la ecuación \eqref{eq:u_sol1} queda expresada por
\begin{equation*}
    -2ub^2\sech^2(bx)+icA_1b\sech^2(bx)+u(1-A_1^2\tanh^2(bx)-\alpha A_2^2\sech^2(bx))=0
\end{equation*}

Esto nos implica que para que exista solución debemos establecer que $c=0$,asi agrupamos los términos y obtenemos que
\begin{equation*}
    u(-2b^2\sech^2(bx)+1-A_1^2\tanh^2(bx)-\alpha A_2^2\sech^2(bx))=0
\end{equation*}

Usamos la identidad $\tanh^2(bx)=1-\sech^2(bx)$ entonces tenemos
\begin{equation*}
    u(-2b^2\sech^2(bx)+1-A_1^2(1-\sech^2(bx))-\alpha A_2^2\sech^2(bx))=0
\end{equation*}
\begin{equation*}
    u[-2b^2\sech^2(bx)+1-A_1^2+A_1^2\sech^2(bx)-\alpha A_2^2\sech^2(bx)]=0
\end{equation*}

Esto nos implica que $A_1=1$ y usando esto obtenemos la siguiente condición $1-2b^2-\alpha A_2^2=0$.

Ahora veamos cómo nos queda la ecuación \eqref{eq:v_sol1}, donde tenemos
\begin{itemize}
    \item $\dot{v}=A_2b(-\tanh(bx)\cdot \sech(bx))=-b\tanh(bx)v$
    \item $\ddot{v}=-b\tanh(bx)\dot{v}-b^2\sech^2(bx)v=b^2\tanh^2(bx)v-b^2\sech^2(bx)v=b^2v(1-2\sech^2(bx))$
\end{itemize}

Así reemplazando estos valores obtenemos
\begin{equation*}
    v(\lambda +b^2-2b^2\sech^2(bx)+1-\alpha u^2-\beta v^2)=0
\end{equation*}
\begin{equation*}
    v(\lambda +b^2-2b^2\sech^2(bx)+1-\alpha\tanh^2(bx)-\beta A_2^2\sech^2(bx))=0
\end{equation*}

Usamos la identidad $\sech$-$\tanh$:
\begin{equation*}
    v(\lambda +b^2-2b^2\sech^2(bx)+1-\alpha[1-\sech^2(bx)]-\beta A_2^2\sech^2(bx))=0
\end{equation*}

Uniendo todo lo anterior obtenemos las siguientes condiciones para la existencia de la solución en función de $\alpha,\beta, A_1,A_2, \lambda, b$ donde los parámetros conocidos son $\alpha,\beta$:
\begin{align*}
A_1&=1\\
b^2&=\alpha-(\lambda+1)\\
2b^2&= \alpha - \beta A_2^2 \\
2b^2&=1-\alpha A_2^2
\end{align*}

Resolvemos esto:
\begin{align*}
    -\beta A_2^2 + \alpha&=1-\alpha A_2^2 \\
\end{align*}

Así obtenemos que
\begin{align*}
    A_1&=1 \\
    A_2^2&=\frac{1-\alpha}{\alpha-\beta} \\
    b^2&=\frac{\alpha^2-\beta}{2(\alpha-\beta)} \\
    \lambda&=\alpha -1-\frac{\alpha^2-\beta}{2(\alpha-\beta)}
\end{align*}

\begin{table}[H]
\centering
\renewcommand{\arraystretch}{1.8}
\resizebox{\textwidth}{!}{%
\begin{tabular}{|l|c|c|c|c|c|}
\hline
\textbf{Condición} & $A_1^2$ & $A_2^2$ & $b^2$ & $\lambda$ & \textbf{Solución} \\
\hline
$\alpha<1,\;\beta<\alpha^2$
& $=1$ & $=\dfrac{1-\alpha}{\alpha-\beta}>0$ & $=\dfrac{\alpha^2-\beta}{2(\alpha-\beta)}>0$
& $=\alpha-1-b^2$
& Ambas componentes \\
\hline
$\alpha>1,\;\beta>\alpha^2$
& $=1$ & $=\dfrac{1-\alpha}{\alpha-\beta}>0$ & $=\dfrac{\alpha^2-\beta}{2(\alpha-\beta)}>0$
& $=\alpha-1-b^2$
& Ambas componentes \\
\hline
$\alpha=1,\;\beta\neq 1$
& $=1$ & $=0$ & $=\dfrac{1}{2}$
& $=-\dfrac{1}{2}$
& Solo $\psi$  \\
\hline
$\beta=\alpha^2,\;\alpha\neq\beta$
& $=1$ & $=\dfrac{1}{\alpha}$ & $=0$
& $=\alpha-1$
& Degenerada ($b=0$) \\
\hline
$\alpha=\beta$
& $=1$ & --- & --- & ---
& No definida (denominador $=0$) \\
\hline
\end{tabular}}
\caption{Casos de existencia de la Solución~1 (ansatz tanh--sech).}
\label{tab:sol1_existencia}
\end{table}

Donde las regiones de solución se representan por 

\begin{figure}[H]
    \centering
    \includegraphics[width=0.5\textwidth]{TESIS/1.1.png}
    \caption{Regiones de existencia de la Solución 1 (ansatz tanh--sech).}
    \label{fig:regiones_sol1}
\end{figure}



\subsection{Solución 1.2}\label{sec:sol112}
Proponemos los siguientes resultados
\begin{align*}
    u&=A_1\tanh(bx)+iB_1\\
    v&=A_2\sech(bx)
\end{align*}
Calculamos las derivadas de $u$ y $v$ para reemplazarlas en las ecuaciones:
\begin{itemize}
    \item $\dot{u}=A_1b\sech^2(bx)$
    \item $\ddot{u}=-2A_1b^2\cdot \sech^2(bx)\tanh(bx)=-2A_1b^2[1-\tanh^2(bx)]\tanh(bx)$
    \item $|u|^2=A_1^2\tanh^2(bx)+B_1^2$
\end{itemize}

Reemplazamos esto en el sistema:
\begin{equation*}
-2A_1b^2[1-\tanh^2(bx)]\tanh(bx)+icA_1b\sech^2(bx)+u(1-|u|^2-\alpha|v|^2)=0
\end{equation*}
\begin{equation*}
-2A_1b^2[1-\tanh^2(bx)]\tanh(bx)+icA_1b[1-\tanh^2(bx)]+u(1-|u|^2-\alpha|v|^2)=0
\end{equation*}
Expandimos los términos y obtenemos
\begin{equation*}
\begin{split}
    &-2A_1b^2[1-\tanh^2(bx)]\tanh(bx)+icA_1b[1-\tanh^2(bx)]+ \\&[A_1\tanh(bx)+iB_1][1-(A_1^2\tanh^2(bx)+B_1^2)-\alpha A_2^2\sech^2(bx)]=0
\end{split}
\end{equation*}

Queremos llegar a una expresión en función de $\tanh(bx)$. Para simplificar el cálculo usaremos el siguiente cambio de variable $T=\tanh(bx)$:
\begin{align*}
-2A_1b^2[1-T^2]T+icA_1b[1-T^2]+[A_1T+iB_1][1-(A_1^2T^2+B_1^2)-\alpha A_2^2(1-T^2)] &=0 \\
 \Leftrightarrow -2A_1b^2T+2A_1b^2T^3+icA_1b-icA_1bT^2 +[A_1T+iB_1](1-A_1^2T^2-B_1^2-\alpha A_2^2+\alpha A_2^2T^2)&=0
\end{align*}

Separamos por parte real e imaginaria y obtenemos el siguiente sistema:
\begin{align*}
    T^3(2A_1b^2-A_1^3+\alpha A_1A_2^2)+T(-2A_1b^2+A_1-A_1B_1^2-\alpha A_1A_2^2)=0 \\
    i[T^2(-A_1^2B_1+\alpha A_2^2B_1-A_1bc)+B_1-B_1^3-\alpha A_2^2B_1+A_1bc]=0
\end{align*}

Así obtenemos las siguientes condiciones:
\begin{align*}
    2A_1b^2-A_1^3+\alpha A_1A_2^2&=0 \\
    -2A_1b^2+A_1-A_1B_1^2-\alpha A_1A_2^2&=0 \\
    -A_1^2B_1+\alpha A_2^2B_1-A_1bc&=0\\
    B_1-B_1^3-\alpha A_2^2B_1+A_1bc&=0
\end{align*}


Usando desarrollos obtenidos en la parte anterior respecto a las derivadas de v obtenemos que debe cumplir 

\begin{align*}
    v(\lambda +b^2-2b^2\sech^2(bx)+1-\alpha(A_1^2\tanh^2(bx)+B_1^2 )-\beta A_2^2\sech^2(bx))&=0 \\
    \Leftrightarrow  v(\lambda +b^2-2b^2\sech^2(bx)+1-\alpha(A_1^2[1-\sech^2(bx)]+B_1^2 )-\beta A_2^2\sech^2(bx))&=0\\
    \Leftrightarrow v(\lambda +b^2-2b^2\sech^2(bx)+1-\alpha A_1^2+\alpha A_1^2\sech^2(bx)-\alpha B_1^2-\beta A_2^2\sech^2(bx))&=0
\end{align*}

Luego las condiciones que se encuentran de esta ecuación son:
\begin{align*}
\lambda+b^2+1-\alpha A_1^2-\alpha B_1^2&=0\\
-2b^2+\alpha A_1^2-\beta A_2^2&=0
\end{align*}


Luego todas las condiciones del sistema corresponden a
\begin{align}
    2A_1b^2-A_1^3+\alpha A_1A_2^2&=0 \label{eq:1.1.2.1}\\
    -2A_1b^2+A_1-A_1B_1^2-\alpha A_1A_2^2&=0 \label{eq:1.1.2.2}\\
    -A_1^2B_1+\alpha A_2^2B_1-A_1bc&=0\label{eq:1.1.2.3}\\
    B_1-B_1^3-\alpha A_2^2B_1+A_1bc&=0 \label{eq:1.1.2.4}\\
    b^2+1+\lambda -\alpha A_1^2-\alpha B_1^2&=0 \label{eq:1.1.2.5}\\
    -2b^2+\alpha A_1^2-\beta A_2^2&=0 \label{eq:1.1.2.6}
\end{align}

Resolviendo el sistema anterior con $c$ dado, todas las incógnitas quedan determinadas en función de $\alpha$, $\beta$ y $c$.
 
\medskip
Dividiendo \eqref{eq:1.1.2.1} y \eqref{eq:1.1.2.2}
por $A_1\neq 0$:
\begin{align*}
    A_1^2 - \alpha A_2^2 = 2b^2 \\
    A_1^2 + B_1^2 = 1.
\end{align*}
 
Usando la ecuaciones $ A_1^2 - \alpha A_2^2 = 2b^2$ y $A_1^2 + B_1^2 = 1$ en \eqref{eq:1.1.2.3}:
\begin{equation*}
    B_1(\alpha A_2^2 - A_1^2) = A_1 bc
    \;\Longrightarrow\;
    -2b^2 B_1 = A_1 bc
    \;\Longrightarrow\;
    B_1 = -\frac{A_1 c}{2b}.
\end{equation*}

 
 Combinando \eqref{eq:1.1.2.6} con $A_1^2 - \alpha A_2^2 = 2b^2$
\begin{equation*}
    \beta A_2^2 = \alpha A_1^2 - 2b^2 = 2b^2(\alpha-1) + \alpha^2 A_2^2,
\end{equation*}
Con lo que obtenemos 
\begin{align*}
A_2^2 &= 2b^2\frac{(\alpha-1)}{(\beta-\alpha^2)} \\
A_1^2 &= 2b^2\frac{(\beta-\alpha)}{(\beta-\alpha^2)}
\end{align*}
\medskip
Sustituyendo $B_1 = \frac{-A_1c}{2b}$ en $A_1^2+B_1^2=1$
se obtiene $A_1^2 = \frac{4b^2}{(4b^2+c^2)}$. Igualando con la expresión $A_1^2 = 2b^2\frac{(\beta-\alpha)}{(\beta-\alpha^2)}$
\begin{equation*}
    \frac{4b^2}{4b^2+c^2} = \frac{2b^2(\beta-\alpha)}{\beta-\alpha^2}
    \;\Longrightarrow\;
    4b^2+c^2 = \frac{2(\beta-\alpha^2)}{\beta-\alpha},
\end{equation*}
y despejando $b^2$:
\begin{equation}
    b^2 = \frac{2(\beta-\alpha^2)-c^2(\beta-\alpha)}{4(\beta-\alpha)}.
    \label{eq:b2final}
\end{equation}
 
\medskip
\noindent\textbf{Valores explícitos en términos de $\alpha$, $\beta$ y $c$.}
Sustituyendo \eqref{eq:b2final} en las expresiones encontradas anteriormente llegamos a los siguientes valores de las variables
 
\begin{align}
    b^2     &= \frac{(\beta-\alpha^2)}{2(\beta-\alpha)}-\frac{c^2}{4}
    \label{eq:bfinal}\\[8pt]
    A_1^2   &= 1-\frac{c^2(\beta-\alpha)}{2(\beta-\alpha^2)}
    \label{eq:A1final}\\[8pt]
    A_2^2   &= \frac{(\alpha-1)\bigl[2(\beta-\alpha^2)-c^2(\beta-\alpha)\bigr]}
                    {2(\beta-\alpha)(\beta-\alpha^2)},
    \label{eq:A2final}\\[8pt]
    B_1^2   &= \frac{c^2(\beta-\alpha)}{2(\beta-\alpha^2)},
    \label{eq:B1final}\\[8pt]
    \lambda &= \alpha - 1 -\frac{(\beta-\alpha^2)}{2(\beta-\alpha)}+\frac{c^2}{4}    
    \label{eq:lambdafinal}
\end{align}
 
\noindent Para $c=0$ se recuperan exactamente los resultados de la Sección~1.1.1.
 
\medskip
\noindent\textbf{Condiciones de existencia.}
Para que la solución sea no trivial ($b^2>0$, $A_1^2>0$, $A_2^2>0$, $B_1^2\geq 0$)
se requiere simultáneamente:
\begin{equation}
    \frac{2(\beta-\alpha^2)}{\beta-\alpha} > c^2, \qquad
    \frac{\alpha-1}{\beta-\alpha^2}>0, \qquad
    \frac{\beta-\alpha}{\beta-\alpha^2}>0.
\end{equation}


\begin{table}[H]
\centering
\renewcommand{\arraystretch}{2.5}
\resizebox{\textwidth}{!}{%
\begin{tabular}{|p{3cm}|c|c|c|c|c|c|}
\hline
\textbf{Condición} & $b^2$ & $A_1^2$ & $A_2^2$ & $B_1^2$ & $\lambda$ & \textbf{Solución} \\
\hline
$\alpha<1,\;\beta<\alpha^2$
\newline
\hspace*{-\tabcolsep}\rule{\dimexpr\linewidth+2\tabcolsep\relax}{0.4pt}
\newline $\beta<\alpha $
\newline
\hspace*{-\tabcolsep}\rule{\dimexpr\linewidth+2\tabcolsep\relax}{0.4pt}
\newline $c^2<\dfrac{2(\beta-\alpha^2)}{\beta-\alpha}$
& $\dfrac{2(\beta-\alpha^2)-c^2(\beta-\alpha)}{4(\beta-\alpha)}>0$
& $\dfrac{2(\beta-\alpha^2)-c^2(\beta-\alpha)}{2(\beta-\alpha^2)}>0$
& $\dfrac{(\alpha-1)[2(\beta-\alpha^2)-c^2(\beta-\alpha)]}{2(\beta-\alpha)(\beta-\alpha^2)}>0$
& $\dfrac{c^2(\beta-\alpha)}{2(\beta-\alpha^2)}>0$
& $\alpha-1-b^2$
& Ambas componentes \\
\hline
$\alpha>1,\;\beta>\alpha^2$
\newline
\hspace*{-\tabcolsep}\rule{\dimexpr\linewidth+2\tabcolsep\relax}{0.4pt}
\newline $\beta>\alpha$
\newline
\hspace*{-\tabcolsep}\rule{\dimexpr\linewidth+2\tabcolsep\relax}{0.4pt}
\newline$c^2<\dfrac{2(\beta-\alpha^2)}{\beta-\alpha}$
& $\dfrac{2(\beta-\alpha^2)-c^2(\beta-\alpha)}{4(\beta-\alpha)}>0$
& $\dfrac{2(\beta-\alpha^2)-c^2(\beta-\alpha)}{2(\beta-\alpha^2)}>0$
& $\dfrac{(\alpha-1)[2(\beta-\alpha^2)-c^2(\beta-\alpha)]}{2(\beta-\alpha)(\beta-\alpha^2)}>0$
& $\dfrac{c^2(\beta-\alpha)}{2(\beta-\alpha^2)}>0$
& $\alpha-1-b^2$
& Ambas componentes \\
\hline
$\alpha=1,\;\beta\neq 1$
\newline
\hspace*{-\tabcolsep}\rule{\dimexpr\linewidth+2\tabcolsep\relax}{0.4pt}
\newline $c^2<2$
& $\dfrac{2-c^2}{4}>0$
& $\dfrac{2-c^2}{2}>0$
& $0$
& $\dfrac{c^2}{2}\geq 0$
& $-\dfrac{2-c^2}{4}$
& Solo $\psi$ \\
\hline
$\beta=\alpha^2$
\newline
\hspace*{-\tabcolsep}\rule{\dimexpr\linewidth+2\tabcolsep\relax}{0.4pt}
\newline$\alpha\neq\beta,\;c\neq 0$
& $-\dfrac{c^2}{4}<0$
& ---
& ---
& ---
& ---
& No existe \\
\hline
$\beta=\alpha^2,\;c=0$
\newline
\hspace*{-\tabcolsep}\rule{\dimexpr\linewidth+2\tabcolsep\relax}{0.4pt}
\newline $\alpha\neq\beta$
& $0$
& $1$
& $\dfrac{1}{\alpha}$
& $0$
& $\alpha-1$
& Degenerada ($b=0$) \\
\hline
$c=0$
& $\dfrac{\beta-\alpha^2}{2(\beta-\alpha)}$
& $1$
& $\dfrac{\alpha-1}{\beta-\alpha}$
& $0$
& $\alpha-1-b^2$
& Solución sin velocidad \\
\hline
$\beta=\alpha$
& ---
& ---
& ---
& ---
& ---
& No definida \\
\hline
\end{tabular}}
\caption{Casos de existencia de la Solución~1.1.2
(ansatz $u=A_1\tanh(bx)+iB_1$, $v=A_2\operatorname{sech}(bx)$, $c\neq 0$).}
\label{tab:sol112_existencia}
\end{table}

\newpage
 Luego nuestra solucion queda escrita como 

 \begin{align*}
\psi(t,x) &=
\sqrt{\frac{2(\beta-\alpha^2)-c^2(\beta-\alpha)}{2(\beta-\alpha^2)}}
\tanh\!\left(\sqrt{\frac{2(\beta-\alpha^2)-c^2(\beta-\alpha)}{4(\beta-\alpha)}}\,(x-ct)\right)
+\,i\sqrt{\frac{c^2(\beta-\alpha)}{2(\beta-\alpha^2)}}
\\[12pt]
\phi(t,x) &=
\sqrt{\frac{(\alpha-1)\bigl[2(\beta-\alpha^2)-c^2(\beta-\alpha)\bigr]}{2(\beta-\alpha)(\beta-\alpha^2)}}
\;\operatorname{sech}\!\left(\sqrt{\frac{2(\beta-\alpha^2)-c^2(\beta-\alpha)}{4(\beta-\alpha)}}\,(x-ct)\right)
\\[6pt]
&\quad\times
\exp\!\left(i\left[\left(\alpha-1-\frac{2(\beta-\alpha^2)-c^2(\beta-\alpha)}{4(\beta-\alpha)}+\frac{c^2}{4}\right)t
-\frac{c}{2}\,x\right]\right)
\end{align*}

\begin{figure}[H]
    \centering
    \includegraphics[width=0.5\textwidth]{TESIS/1.2.1.png}
    \caption{Regiones de existencia de la Solución 1 (ansatz tanh--sech).}
    \label{fig:regiones_sol1}
\end{figure}

\begin{figure}[H]
    \centering
    \includegraphics[width=0.5\textwidth]{TESIS/1.2.2.png}
    \caption{Regiones de existencia de la Solución 1 (ansatz tanh--sech).}
    \label{fig:regiones_sol1}
\end{figure}


\begin{figure}[H]
    \centering
    \includegraphics[width=0.5\textwidth]{TESIS/1.2.3.png}
    \caption{Regiones de existencia de la Solución 1 (ansatz tanh--sech).}
    \label{fig:regiones_sol1}
\end{figure}

\subsection{Solución 1.3 --- Caso desacoplado ($\phi=0$)}

Cuando $\phi = 0$ la ecuación de $\phi$ se satisface trivialmente y el sistema se reduce a
\[
  ic\dot{u} + \ddot{u} + u(1 - |u|^2) = 0.
\]
Proponemos como solucion

\begin{equation*}
    u = A_1\tanh(bx) + iB_1
\end{equation*}

Los cálculos son idénticos a partes anteriores donde tenemos 
\begin{equation*}
\begin{split}
    &-2A_1b^2[1-\tanh^2(bx)]\tanh(bx)+icA_1b[1-\tanh^2(bx)]+ \\&[A_1\tanh(bx)+iB_1][1-(A_1^2\tanh^2(bx)+B_1^2)]=0
\end{split}
\end{equation*}


Asi nuestro sistema nos queda de la forma 

\begin{align*}
    2A_1b^2-A_1^3&=0 \\
    -2A_1b^2+A_1-A_1B_1^2&=0 \\
    -A_1^2B_1-A_1bc&=0\\
    B_1-B_1^3+A_1bc&=0
\end{align*}

Simplificamos algunas ecuaciones por $A_1$
\begin{align}
    2b^2-A_1^2&=0 \label{3.1} \\
    -2b^2+1-B_1^2&=0 \label{3.2}\\
    -A_1B_1-bc&=0\\
    B_1-B_1^3+A_1bc&=0
\end{align}

Sumamos las ecuaciones \eqref{3.1} y \eqref{3.2} donde obtenemos

\begin{equation*}
    A_1^2+B_1^2=1
\end{equation*}

Ademas obtenemos que 

\begin{align}
    B_1^2&=\frac{c^2}{2} \\
    b^2&=\frac{1}{2}-\frac{c^2}{4} \\
    A_1^2&=1-\frac{c^2}{2}
\end{align}

Así notamos que la condición de existencia es $\sqrt{2}\geq |c|$. \textbf{Cuando c=0} tenemos que 

 $b = \frac{\sqrt{2}}{2}$, $A_1 = 1$, $B_1 = 0$



\begin{table}[H]
\centering
\renewcommand{\arraystretch}{1.8}
\begin{tabular}{|l|c|c|c|}
\hline
\textbf{Condición} & $A_1^2$ & $B_1^2$ & $b^2$ \\
\hline
$|c| < \sqrt{2}$ & $>0$ & $>0$ & $>0$ \\
\hline
$c = 0$ & $=1$ & $=0$ & $=\tfrac{1}{2}$ \\
\hline
$|c| = \sqrt{2}$ & $=0$ & $=1$ & $=0$ \\
\hline
\end{tabular}
\caption{Casos de existencia de la Solución 1.3 (desacoplada).}
\label{tab:122}
\end{table}

\begin{figure}[H]
    \centering
    \includegraphics[width=0.8\textwidth]{TESIS/img/Solcuion1.3.png}
   
    \label{fig:regiones_sol3}
\end{figure}

Así nuestra solución explicita es 

\begin{align*}
    \psi(y)=\sqrt{\frac{2-c^2}{2}}\tanh{\left(\frac{\sqrt{2-c^2}}{2}y\right)}+i\frac{c}{\sqrt{2}}
\end{align*}

\subsection{Solución 1.4}

Proponemos los siguientes resultados
\begin{align*}
     u = A_1\tanh(bx) + iB_1 \\
     v = A_2\tanh(bx)
\end{align*}



A diferencia de la Solución 1.2 aquí $v$ tiene perfil tipo $\tanh$ en lugar de $\sech$. Usando esta solución y $T=\tanh(bx)$ obtenemos que 
\begin{equation}
    -2A_1b^2[1-T^2]T+icA_1b[1-T^2]+[A_1T+iB_1][1-(A_1^2T^2+B_1^2)-\alpha A_2^2T^2] =0
\end{equation}


\emph{Parte real:}
\begin{align}
  1 - B_1^2 &= 2b^2, \label{eq:123r1}\\
  A_1^2 + \alpha A_2^2 &= 2b^2. \label{eq:123r2}
\end{align}

\emph{Parte imaginaria:}
\begin{align}
    cA_1b+B_1-B_1^3=0 \\
    cA_1b+B_1(A_1^2+\alpha A_2^2)=0
\end{align}
\textbf{Ecuación de $v$.}
Sustituyendo en la ecuación~\eqref{eq:v_sol1} obtenemos las siguientes ecuaciones 
\begin{align}
  1 + \lambda - \alpha B_1^2 &= 2b^2, \label{eq:123v1}\\
  \alpha A_1^2 + \beta A_2^2 &= 2b^2. \label{eq:123v2}
\end{align}


Con esto obtenemos el sistema 
\begin{align}
  1 - B_1^2 &= 2b^2 \label{eq:4.1}\\
  A_1^2 + \alpha A_2^2 &= 2b^2  \label{eq:4.2} \\
  cA_1b+B_1-B_1^3&=0 \label{eq:4.3}  \\
  cA_1b+B_1(A_1^2+\alpha A_2^2)&=0 \label{eq:4.4} \\
  1 + \lambda - \alpha B_1^2 &= 2b^2 \label{eq:4.5} \\
  \alpha A_1^2 + \beta A_2^2 &= 2b^2 \label{eq:4.6}
\end{align}

\textbf{Resolución.}
De las ecuaciones~\eqref{eq:4.2} y~\eqref{eq:4.6} se tiene
$A_1^2 + \alpha A_2^2 = \alpha A_1^2 + \beta A_2^2$, luego
\[
  A_1^2= \frac{\beta - \alpha}{1 - \alpha} A_2^2
\]
y sustituyendo en~\eqref{eq:4.2}
\[
  A_1^2 = \frac{2b^2(\beta-\alpha)}{\beta-\alpha^2},
  \qquad
  A_2^2 = \frac{2b^2(1-\alpha)}{\beta-\alpha^2}.
\]


Sustituyendo \eqref{eq:4.1} en \eqref{eq:4.3} y la expresión de $A_1^2$ obtenemos que 

\begin{equation}
    B_1^2=\frac{c^2(\beta-\alpha)}{2(\beta-\alpha^2)}
\end{equation}

Sustituyendo \eqref{eq:4.1} obtenemos que 

\begin{equation}
    b^2=\frac{2(\beta-\alpha^2)-c^2(\beta-\alpha)}{4(\beta-\alpha^2)}
\end{equation}

Denotamos $N=2(\beta-\alpha^2)-c^2(\beta-\alpha)$ entonces

\begin{align*}
    A_1^2&=\frac{N(\beta-\alpha)}{2(\beta-\alpha^2)^2} \\
    A_2^2&=\frac{N(1-\alpha)}{2(\beta-\alpha^2)^2} \\
    B_1^2&=\frac{c^2(\beta-\alpha)}{2(\beta-\alpha^2)} \\
    \lambda&=\frac{(\alpha-1)(\beta-\alpha)c^2}{2(\beta-\alpha^2)}
\end{align*}

\begin{table}[H]
\centering
\renewcommand{\arraystretch}{1.8}
\resizebox{\textwidth}{!}{%
\begin{tabular}{|p{3cm}|c|c|c|c|c|c|}
\hline
\textbf{Condición} & $b^2$ & $A_1^2$ & $A_2^2$ & $B_1^2$ & $\lambda$ & \textbf{Solución} \\
\hline
$\alpha<1,\;\beta>\alpha$
\newline
\hspace*{-\tabcolsep}\rule{\dimexpr\linewidth+2\tabcolsep\relax}{0.4pt}
\newline $\beta>\alpha^2$
\newline
\hspace*{-\tabcolsep}\rule{\dimexpr\linewidth+2\tabcolsep\relax}{0.4pt}
\newline $c^2<\dfrac{2(\beta-\alpha^2)}{\beta-\alpha}$
& $\dfrac{\mathcal{N}}{4(\beta-\alpha^2)}>0$
& $\dfrac{\mathcal{N}(\beta-\alpha)}{2(\beta-\alpha^2)^2}>0$
& $\dfrac{\mathcal{N}(1-\alpha)}{2(\beta-\alpha^2)^2}>0$
& $\dfrac{c^2(\beta-\alpha)}{2(\beta-\alpha^2)}>0$
& $\dfrac{(\alpha-1)(\beta-\alpha)c^2}{2(\beta-\alpha^2)}$
& Ambas componentes \\
\hline
$\alpha>1,\;\beta<\alpha$
\newline
\hspace*{-\tabcolsep}\rule{\dimexpr\linewidth+2\tabcolsep\relax}{0.4pt}
\newline $\beta<\alpha^2$
\newline
\hspace*{-\tabcolsep}\rule{\dimexpr\linewidth+2\tabcolsep\relax}{0.4pt}
\newline $c^2<\dfrac{2(\beta-\alpha^2)}{\beta-\alpha}$
& $\dfrac{\mathcal{N}}{4(\beta-\alpha^2)}>0$
& $\dfrac{\mathcal{N}(\beta-\alpha)}{2(\beta-\alpha^2)^2}>0$
& $\dfrac{\mathcal{N}(1-\alpha)}{2(\beta-\alpha^2)^2}>0$
& $\dfrac{c^2(\beta-\alpha)}{2(\beta-\alpha^2)}>0$
& $\dfrac{(\alpha-1)(\beta-\alpha)c^2}{2(\beta-\alpha^2)}$
& Ambas componentes \\
\hline
$\alpha=1,\;\beta\neq 1$
\newline
\hspace*{-\tabcolsep}\rule{\dimexpr\linewidth+2\tabcolsep\relax}{0.4pt}
\newline $c^2<2$
& $\dfrac{2-c^2}{4}>0$
& $\dfrac{2-c^2}{2(\beta-1)}>0$
& $0$
& $\dfrac{c^2}{2}\geq 0$
& $0$
& Solo $\psi$ \\
\hline
$\beta=\alpha^2$
\newline
\hspace*{-\tabcolsep}\rule{\dimexpr\linewidth+2\tabcolsep\relax}{0.4pt}
\newline $\alpha\neq\beta,\;c\neq 0$
& $-\dfrac{c^2}{4}<0$
& ---
& ---
& ---
& ---
& No existe \\
\hline
$c=0$
& $\dfrac{1}{2}$
& $\dfrac{\beta-\alpha}{\beta-\alpha^2}$
& $\dfrac{1-\alpha}{\beta-\alpha^2}$
& $0$
& $0$
& Solucion sin velocidad \\
\hline
$\beta=\alpha$
& ---
& ---
& ---
& ---
& ---
& No definida \\
\hline
\end{tabular}}
\caption{Casos de existencia para el ansatz
$u=A_1\tanh(bx)+iB_1$, $v=A_2\tanh(bx)$, $c\neq 0$.
Se denota $\mathcal{N}=2(\beta-\alpha^2)-c^2(\beta-\alpha)$.}
\label{tab:sol_tanh_tanh_cneq0}
\end{table}

La solucion en forma explicita esta representada por 

\begin{align*}
\psi(t,x) &=
\sqrt{\frac{\mathcal{N}\,(\beta-\alpha)}{2(\beta-\alpha^2)^2}}
\tanh\!\left(\sqrt{\frac{\mathcal{N}}{4(\beta-\alpha^2)}}\,(x-ct)\right)
+\,i\,\frac{c\sqrt{\beta-\alpha}}{\sqrt{2(\beta-\alpha^2)}}
\\[12pt]
\phi(t,x) &=
\sqrt{\frac{\mathcal{N}\,(1-\alpha)}{2(\beta-\alpha^2)^2}}
\;\tanh\!\left(\sqrt{\frac{\mathcal{N}}{4(\beta-\alpha^2)}}\,(x-ct)\right)
\\[6pt]
&\quad\times
\exp\!\left(i\left[\left(\frac{(\alpha-1)(\beta-\alpha)c^2}{2(\beta-\alpha^2)}
+\frac{c^2}{4}\right)t
-\frac{c}{2}\,x\right]\right)
\end{align*}
donde $\mathcal{N} = 2(\beta-\alpha^2)-c^2(\beta-\alpha)$.

\begin{figure}[H]
    \centering
    \includegraphics[width=0.5\textwidth]{TESIS/1.4.1.png}
    \caption{Regiones de existencia de la Solución 1 (ansatz tanh--stanh).}
    \label{fig:regiones_sol1}
\end{figure}

\begin{figure}[H]
    \centering
    \includegraphics[width=0.5\textwidth]{TESIS/1.4.2.png}
    \caption{Regiones de existencia de la Solución 1 (ansatz tanh--tanh).}
    \label{fig:regiones_sol1}
\end{figure}

\section{Solución 2}
Ahora utilizaremos ondas viajeras de la siguiente forma
\begin{align*}
    \psi(t,x)&=u(x-ct)\\
    \phi(t,x)&=v(x-ct)
\end{align*}
Donde obtenemos que el sistema nos queda de la siguiente forma
 \begin{align}
     ic\dot{u}+\ddot{u}&+u(1-|u|^2-\alpha|v|^2)=0 \label{eq:u_sol2}\\
    ic\dot{v}+\ddot{v}&+v(1-\alpha|u|^2-\beta|v|^2)=0 \label{eq:v_sol2}
\end{align}

Luego estudiaremos soluciones de distintos tipos.

\subsection{Solución 2.1}\label{sec:sol21}
Establecemos que $c=0$ donde las soluciones son de la forma
\begin{align*}
    \psi(t,x)&=u(x)\\
    \phi(t,x)&=v(x)
\end{align*}

\subsubsection{Solución 2.1.1}\label{sec:sol211}

Proponemos los siguientes resultados
\[
  u = A_1 \tanh(bx), \qquad v = A_2 \tanh(bx).
\]

Calculamos las derivadas de $u$ y $v$. El cálculo es simétrico para ambas; presentamos el de $u$:
\[
  \dot{u}  = A_1 b\,\sech^2(bx),
  \qquad
  \ddot{u} = -2A_1 b^2\sech^2(bx)\tanh(bx).
\]

Introducimos el cambio de variable $T = \tanh(bx)$ y usamos $\sech^2(bx) = 1 - T^2$.
Sustituyendo en la ecuación de $u$, con $|u|^2 = A_1^2 T^2$ y $|v|^2 = A_2^2 T^2$:
\[
  -2A_1 b^2(1-T^2)T + A_1 T\bigl(1 - A_1^2 T^2 - \alpha A_2^2 T^2\bigr) = 0.
\]
Factorizamos $A_1 T$ y reagrupamos por potencias de $T$:
\[
  A_1 T\Bigl[(1 - 2b^2) + T^2\bigl(2b^2 - A_1^2 - \alpha A_2^2\bigr)\Bigr] = 0.
\]
Anulando cada coeficiente obtenemos las condiciones
\begin{align}
  1 - 2b^2 &= 0, \label{eq:211cond1}\\
  A_1^2 + \alpha A_2^2 &= 2b^2. \label{eq:211cond2}
\end{align}

El procedimiento es completamente análogo para $v$, dando lugar a
\begin{align}
  1 - 2b^2 &= 0, \label{eq:211cond3}\\
  \alpha A_1^2 + \beta A_2^2 &= 2b^2. \label{eq:211cond4}
\end{align}

La condición~\eqref{eq:211cond1} determina inmediatamente el parámetro $b$:
\[
  \boxed{b^2 = \tfrac{1}{2}.}
\]

Nótese que $b$ queda \emph{completamente fijado} por el sistema, sin ser un parámetro libre.
Sustituyendo $2b^2 = 1$ en~\eqref{eq:211cond2} y~\eqref{eq:211cond4} obtenemos el sistema lineal
\begin{align*}
  A_1^2 + \alpha A_2^2 &= 1,\\
  \alpha A_1^2 + \beta A_2^2 &= 1.
\end{align*}

Multiplicando la primera ecuación por $\alpha$ y restando la segunda:
\[
  (\alpha^2 - \beta)\,A_2^2 = \alpha - 1,
\]
de donde
\[
  A_2^2 = \frac{1 - \alpha}{\beta - \alpha^2},
  \qquad
  A_1^2 = 1 - \alpha A_2^2 = \frac{\beta - \alpha}{\beta - \alpha^2}.
\]

Las condiciones de existencia de la solución son $\beta > \alpha^2$, $\beta > \alpha$
(para $A_1^2 > 0$) y $\alpha < 1$ (para $A_2^2 > 0$).
Así las soluciones de la Solución~\ref{sec:sol21} son de la forma
\[
  \psi(t,x) = \sqrt{\frac{\beta-\alpha}{\beta-\alpha^2}}\;\tanh\!\left(\frac{x}{\sqrt{2}}\right),
  \qquad
  \phi(t,x) = \sqrt{\frac{1-\alpha}{\beta-\alpha^2}}\;\tanh\!\left(\frac{x}{\sqrt{2}}\right).
\]

\subsection{Solución 2.2}\label{sec:sol22}

Establecemos que $c \neq 0$. El sistema reducido es
\begin{align}
  ic\dot{u} + \ddot{u} + u(1 - |u|^2 - \alpha|v|^2) &= 0, \label{eq:22sist1}\\
  ic\dot{v} + \ddot{v} + v(1 - \alpha|u|^2 - \beta|v|^2) &= 0. \label{eq:22sist2}
\end{align}

\paragraph{Incompatibilidad con perfiles reales.}
Si tomamos $u$ y $v$ reales, el término $ic\dot{u}$ es imaginario puro mientras que el resto
de la ecuación~\eqref{eq:22sist1} es real. Por independencia de partes real e imaginaria se
requeriría $c\dot{u} = 0$, lo que implica $c = 0$. El mismo argumento se aplica a la ecuación de $v$. Por lo tanto, ningún ansatz real es compatible con $c \neq 0$, en completa analogía con lo mostrado en la Sección~\ref{sec:sol32}.

\paragraph{Ansatz compatible.}
La presencia simultánea de los términos $ic\dot{u}$ e $ic\dot{v}$ obliga a que \emph{ambos}
perfiles sean complejos. Proponemos
\[
  u = A_1\tanh(bx) + iB_1,
  \qquad
  v = A_2\tanh(bx) + iD_2,
\]
con $A_1, A_2, B_1, D_2 \in \mathbb{R}$. Entonces
\[
  \dot{u} = A_1 b\,\sech^2(bx), \qquad
  \dot{v} = A_2 b\,\sech^2(bx),
\]
y los módulos cuadrados son
\[
  |u|^2 = A_1^2 T^2 + B_1^2, \qquad
  |v|^2 = A_2^2 T^2 + D_2^2, \qquad T = \tanh(bx).
\]

\paragraph{Ecuación de $u$.}
Separando partes real e imaginaria de~\eqref{eq:22sist1} con $S^2 = 1 - T^2$:

\emph{Parte real:}
\[
  A_1 T\Bigl[(1 - B_1^2 - \alpha D_2^2 - 2b^2) + T^2(2b^2 - A_1^2 - \alpha A_2^2)\Bigr] = 0,
\]
de donde
\begin{align}
  1 - B_1^2 - \alpha D_2^2 &= 2b^2, \label{eq:22r1}\\
  A_1^2 + \alpha A_2^2 &= 2b^2. \label{eq:22r2}
\end{align}

\emph{Parte imaginaria} (usando $icA_1 b S^2 = icA_1 b(1-T^2)$ e igualando
coeficientes de la constante y de $T^2$):
\begin{equation}
  A_1\,c = -2b\,B_1. \label{eq:22i1}
\end{equation}

\paragraph{Ecuación de $v$.}
El cálculo es completamente análogo y produce
\begin{align}
  1 - \alpha B_1^2 - \beta D_2^2 &= 2b^2, \label{eq:22r3}\\
  \alpha A_1^2 + \beta A_2^2 &= 2b^2, \label{eq:22r4}\\
  A_2\,c &= -2b\,D_2. \label{eq:22i2}
\end{align}

\paragraph{Resolución.}
De las ecuaciones~\eqref{eq:22r2} y~\eqref{eq:22r4} se tiene
$A_1^2 + \alpha A_2^2 = \alpha A_1^2 + \beta A_2^2$, luego
\[
  \frac{A_1^2}{A_2^2} = \frac{\beta - \alpha}{1 - \alpha}.
\]
Sustituyendo en~\eqref{eq:22r2}:
\[
  A_2^2\,\frac{\beta - \alpha^2}{1-\alpha} = 2b^2,
\]
de donde
\[
  A_1^2 = \frac{2b^2(\beta-\alpha)}{\beta-\alpha^2},
  \qquad
  A_2^2 = \frac{2b^2(1-\alpha)}{\beta-\alpha^2}.
\]

De las relaciones~\eqref{eq:22i1} e~\eqref{eq:22i2}:
\[
  B_1^2 = \frac{A_1^2 c^2}{4b^2} = \frac{(\beta-\alpha)\,c^2}{2(\beta-\alpha^2)},
  \qquad
  D_2^2 = \frac{A_2^2 c^2}{4b^2} = \frac{(1-\alpha)\,c^2}{2(\beta-\alpha^2)}.
\]

Sustituyendo en~\eqref{eq:22r1}:
\[
  1 - \frac{(\beta-\alpha)\,c^2}{2(\beta-\alpha^2)} - \alpha\cdot\frac{(1-\alpha)\,c^2}{2(\beta-\alpha^2)}
  = 2b^2.
\]
El numerador del segundo y tercer término se simplifica:
\[
  (\beta-\alpha) + \alpha(1-\alpha) = \beta - \alpha^2,
\]
de modo que
\[
  1 - \frac{(\beta-\alpha^2)\,c^2}{2(\beta-\alpha^2)} = 2b^2,
\]
y se obtiene el resultado notable
\[
  \boxed{b^2 = \frac{1}{2} - \frac{c^2}{4}.}
\]

El parámetro $b$ queda completamente determinado por $c$ y es independiente de $\alpha$ y $\beta$.
La condición de existencia $b^2 > 0$ requiere $|c| < \sqrt{2}$.
Cuando $c = 0$ se recupera $b^2 = \tfrac{1}{2}$ y $B_1 = D_2 = 0$, reduciendo la solución a la
Solución~\ref{sec:sol211}.


\section{Solución 3}

Consideramos ondas viajeras con fases independientes para cada componente:
\begin{align*}
    \psi(t,x)&=u(x-ct)e^{i\lambda_1t}\\
    \phi(t,x)&=v(x-ct)e^{i\lambda_2t}
    \label{eq:ansatz3}
\end{align*}
donde $u,v:\mathbb{R}\to\mathbb{R}$ son funciones reales, $c\in\mathbb{R}$ es la velocidad de la onda
y $\lambda_1,\lambda_2\in\mathbb{R}$ son frecuencias a determinar.

\medskip
\noindent\textbf{Derivación del sistema reducido.}
Calculamos las derivadas parciales. Para $\psi$:
\begin{align*}
    i\partial_t\psi &= (-ic\dot{u} - \lambda_1 u)\,e^{i\lambda_1 t} \\
    \partial_{xx}\psi &= \ddot{u}\,e^{i\lambda_1 t}
\end{align*}

Lo cual es analago para $\phi$, asi sustituyendo en las ecuaciones \eqref{eq:NLS1} y \eqref{eq:NLS2}, cancelando los factores de fase que no se anulan $e^{i\lambda_1 t}\neq 0$ y $e^{i\lambda_2 t}\neq 0$ y usando que $|\psi|^2=u^2$, $|\phi|^2=v^2$ (con $u,v$ reales) es posible obtener


\begin{align}
    -ic\dot{u}-\lambda_1u&=\ddot{u}+u(1-u^2-\alpha v^2) \label{eq:u_sol3.1}\\
    -ic\dot{v}-\lambda_2v&=\ddot{v}+v(1-\alpha u^2-\beta v^2) \label{eq:v_sol3.2}
\end{align}

\medskip
\noindent\textbf{Incompatibilidad con $c\neq 0$ para perfiles reales.}
Como $u$ y $v$ son funciones reales, los términos $\ddot{u}$, $u(1+\lambda_1 - u^2-\alpha v^2)$
son reales, mientras que $-ic\dot{u}$ es puramente imaginario. Separando partes real e imaginaria
en el sistema anterior obtenemos  la condicion 
\begin{equation}
    c\dot{u} = 0 \qquad \forall\,x.
    \label{eq:cdu}
\end{equation}
Para un perfil no constante ($\dot{u}\not\equiv 0$) esto fuerza $c=0$. El mismo argumento
aplicado a \eqref{eq:sol3v.2} exige $c\dot{v}=0$. Por lo tanto, \emph{ningún ansatz real
es compatible con $c\neq 0$}, y la velocidad queda determinada por el perfil:
\begin{equation}
    c = 0.
    \label{eq:c0}
\end{equation}
Con $c=0$ las ecuaciones~\eqref{eq:u_sol3.1}--\eqref{eq:v_sol3.2} se reducen a
\begin{align}
    \ddot{u} + u(1+\lambda_1 - u^2 - \alpha v^2) &= 0, \label{eq:sol3u0}\\
    \ddot{v} + v(1+\lambda_2 - \alpha u^2 - \beta v^2) &= 0. \label{eq:sol3v0}
\end{align}
A continuación resolvemos este sistema para distintos ansätze.

%------------------------------------------------------------
\subsection{Solución 3.1 — Ansatz tanh--tanh}
\label{sec:sol31}

Proponemos
\begin{align*}
    u(x) = A_1\tanh(bx) \\
    \qquad v(x) = A_2\tanh(bx)
\end{align*}
Como se observó en~\eqref{eq:cdu}, para este perfil $\dot{u} = A_1 b\operatorname{sech}^2(bx)\neq 0$,
de modo que $c=0$ es la única opción compatible.

Hacemos una verificación rápida en $x=0$, $\sech(0)=1$ y $\sech''(0)=-1$; con nuestra fórmula, $b^2 \cdot 1 \cdot (1-2)=-b^2$. 

El cálculo es simétrico para $v$. Asi tenemos que
\begin{align*}
    b^2u(1-2\sech^2(bx))+u(1+\lambda_1-A_1^2\sech^2(bx)-\alpha A_2^2\sech^2(bx))&=0\\
    u\bigl[(b^2+1+\lambda_1)+(-2b^2-A_1^2-\alpha A_2^2)\sech^2(bx)\bigr]&=0
\end{align*}

Anulando los coeficientes:
\begin{align*}
    b^2+1+\lambda_1&=0, \\
    2b^2+A_1^2+\alpha A_2^2&=0.
\end{align*}

La segunda condición exige $A_1^2+\alpha A_2^2=-2b^2$, lo que es imposible para amplitudes reales cuando $\alpha\geq 0$ (ya que el lado izquierdo es no negativo y el derecho es estrictamente negativo).
Un análisis análogo sobre la ecuación de $v$ produce la condición $\alpha A_1^2+\beta A_2^2=-2b^2$, que es igualmente imposible para $\alpha,\beta\geq 0$.

\medskip
\textbf{Conclusión:} El ansatz $u=A_1\sech(bx)$, $v=A_2\sech(bx)$ con fases $e^{i\lambda_1 t}$, $e^{i\lambda_2 t}$
es \textbf{incompatible} con el sistema \eqref{eq:NLS1}--\eqref{eq:NLS2} para $\alpha,\beta>0$.
Intuitivamente, esto refleja que la ecuación NLS con el signo $+\psi(1-|\psi|^2)$ es de tipo defocalizante,
y los solitones brillantes (perfil $\sech$) requieren no linealidad focalizante.

%------------------------------------------------------------
\subsection{Solución 3.2 — Ansatz sech--sech (imposible)}
\label{sec:sol32}

Proponemos 
\begin{align*}
    u = A_1\operatorname{sech}(bx) \\
    v = A_2\operatorname{sech}(bx)
\end{align*}
Usando
$\ddot{u} = b^2 u(1-2\operatorname{sech}^2(bx))$ (véase la verificación en la sección anterior),
la ecuación~\eqref{eq:sol3u0} produce las condiciones
\begin{equation*}
    b^2 + 1 + \lambda_1 = 0, \qquad
    2b^2 + A_1^2 + \alpha A_2^2 = 0.
\end{equation*}
La segunda condición exige $A_1^2 + \alpha A_2^2 = -2b^2$, lo que es imposible para amplitudes
reales cuando $\alpha\geq 0$, pues el lado izquierdo es no negativo y el derecho estrictamente
negativo. Un análisis análogo sobre la ecuación de $v$ reproduce la misma contradicción.

\medskip
\noindent\textbf{Conclusión.} El ansatz sech--sech es incompatible con el sistema para
$\alpha,\beta>0$, independientemente del valor de $c$.

%------------------------------------------------------------
\subsection{Solución 3.3 — Ansatz sech--tanh}
\label{sec:sol33}

Proponemos

\begin{align*}
     u(x) = A_1\operatorname{sech}(bx) \\
     v(x) = A_2\tanh(bx)
\end{align*}

Para $u$: $\dot{u} = -A_1 b\operatorname{sech}(bx)\tanh(bx)$ (impar), mientras que los demás
términos  son pares. Para $v$: $\dot{v} = A_2 b\operatorname{sech}^2(bx)$
(par), mientras que los demás términos son impares. En ambos casos la condición~\eqref{eq:cdu}
exige $c=0$.

\medskip
\noindent\textbf{Sustitución.} Con $c=0$, usando
$\ddot{u} = b^2u(1-2\operatorname{sech}^2(bx))$,
$\ddot{v} = -2b^2v(1-\tanh^2(bx))$
y la identidad $\tanh^2 = 1-\operatorname{sech}^2$, obtenemos el sistema
\begin{align}
    b^2 + 1 + \lambda_1 - \alpha A_2^2 &= 0, \label{eq:33l1}\\
    2b^2 + A_1^2 - \alpha A_2^2 &= 0, \label{eq:33A1}\\
    -2b^2 + 1 + \lambda_2 - \alpha A_1^2 &= 0, \label{eq:33l2}\\
    2b^2 + \alpha A_1^2 - \beta A_2^2 &= 0. \label{eq:33A2}
\end{align}

\noindent\textbf{Resolución.} De~\eqref{eq:33A1} y~\eqref{eq:33A2}:
\begin{equation}
    A_1^2 = \frac{-2b^2(\beta-\alpha)}{\beta-\alpha^2}, \qquad
    A_2^2 = \frac{2b^2(1-\alpha)}{\beta-\alpha^2}.
    \label{eq:33amps}
\end{equation}
De~\eqref{eq:33l1} y~\eqref{eq:33l2}:
\begin{align}
    \lambda_1 &= -b^2\,\frac{\alpha^2-2\alpha+\beta}{\beta-\alpha^2} - 1, \label{eq:33lam1}\\
    \lambda_2 &= \frac{2b^2\beta(1-\alpha)}{\beta-\alpha^2} - 1. \label{eq:33lam2}
\end{align}
Nótese que en general $\lambda_1\neq\lambda_2$. El parámetro $b>0$ es libre. \textbf{Condiciones de existencia.} Para $A_1^2>0$ se requiere
$(\beta-\alpha)/(\beta-\alpha^2)<0$; para $A_2^2>0$ se requiere
$(1-\alpha)/(\beta-\alpha^2)>0$.

\begin{table}[H]
\centering
\renewcommand{\arraystretch}{1.8}
\begin{tabular}{|p{3.5cm}|c|c|c|}
\hline
\textbf{Condición} & $A_1^2$ & $A_2^2$ & \textbf{Solución} \\
\hline
$\alpha<1,\;\alpha^2<\beta<\alpha$
& $>0$
& $>0$
& Ambas componentes \\
\hline
$\alpha>1,\;\alpha<\beta<\alpha^2$
& $>0$
& $>0$
& Ambas componentes \\
\hline
$\beta=\alpha,\;\beta>\alpha^2$
& $0$
& $>0$
& Solo $\phi$ \\
\hline
Resto & --- & --- & No existe \\
\hline
\end{tabular}
\caption{Casos de existencia de la Solución~3.3 (ansatz sech--tanh, $c=0$).
El parámetro $b>0$ es libre.}
\label{tab:sol33}
\end{table}

\noindent La solución explícita es
\begin{equation*}
    \psi(t,x) = \sqrt{\frac{-2b^2(\beta-\alpha)}{\beta-\alpha^2}}\,
    \operatorname{sech}(bx)\,e^{i\lambda_1 t}, \qquad
    \phi(t,x) = \sqrt{\frac{2b^2(1-\alpha)}{\beta-\alpha^2}}\,
    \tanh(bx)\,e^{i\lambda_2 t},
\end{equation*}
con $\lambda_1$, $\lambda_2$ dados por~\eqref{eq:33lam1}--\eqref{eq:33lam2}
y $b>0$ fijado externamente.

\section{Solucion 4}
Proponemos la solucion 

\begin{align*}
    \psi&=e^{i[(\lambda_1+\frac{c^2}{4})t-c\frac{x}{2}]}u(x-ct) \\
    \phi&=e^{i[(\lambda_2+\frac{c^2}{4})t-c\frac{x}{2}]}v(x-ct) 
\end{align*}

Luego reemplazamos esto en nuestro sistema. Usando lo encontrado para la segunda ecuacion de la solución 1, tenemos que
\begin{align*}
       \ddot{u}&+u (1+\lambda_1-|u|^2-\alpha|v|^2)=0 \\ 
       \ddot{v}&+v (1+\lambda_2-\alpha |u|^2-\beta|v|^2)=0 
\end{align*}
Donde proponemos que 

\begin{align*}
    u(x) = A_1\sech(b_1x) \\
    \qquad v(x) = A_2\sech(b_2x)
\end{align*}

Entonces usando calculos anteriores sabemos que 
\begin{align*}
 b_1^2u(1-2\sech^2(b_1x))+A_1\sech(b_1x)(1+\lambda_1-A_1^2\sech^2(b_1x)-\alpha A_2^2\sech^2(b_2x))=0 \\
b_2^2v(1-2\sech^2(b_2x))+A_2\sech(b_2x)(1+\lambda_2-\alpha A_1^2\sech^2(b_1x)-\beta A_2^2\sech^2(b_2x))=0 
\end{align*}

Donde esto es equivalente a 
\begin{align*}
 b_1^2A_1\sech(b_1x)(1-2\sech^2(b_1x))+A_1\sech(b_1x)(1+\lambda_1-A_1^2\sech^2(b_1x)-\alpha A_2^2\sech^2(b_2x))=0 \\
b_2^2A_2\sech(b_2x)(1-2\sech^2(b_2x))+A_2\sech(b_2x)(1+\lambda_2-\alpha A_1^2\sech^2(b_1x)-\beta A_2^2\sech^2(b_2x))=0 
\end{align*}

Notemos que al dividir por $A_1\sech(b_1 x)$ en la primera ecuación y por 
$A_2\sech(b_2 x)$ en la segunda, obtenemos
\begin{align*}
    b_1^2(1-2\sech^2(b_1x)) + (1+\lambda_1 - A_1^2\sech^2(b_1x) - \alpha A_2^2\sech^2(b_2x)) &= 0 \\
    b_2^2(1-2\sech^2(b_2x)) + (1+\lambda_2 - \alpha A_1^2\sech^2(b_1x) - \beta A_2^2\sech^2(b_2x)) &= 0
\end{align*}
Reordenando, agrupamos los términos constantes y los términos con funciones 
hiperbólicas:
\begin{align*}
    (b_1^2 + 1+\lambda_1) - (2b_1^2 + A_1^2)\sech^2(b_1x) - \alpha A_2^2\sech^2(b_2x) &= 0 \\
    (b_2^2 + 1+\lambda_2) - \alpha A_1^2\sech^2(b_1x) - (2b_2^2 + \beta A_2^2)\sech^2(b_2x) &= 0
\end{align*}
Para que estas identidades se satisfagan para todo $x\in\mathbb{R}$, cada 
función linealmente independiente debe tener coeficiente nulo. Sin embargo, 
cuando $b_1 \neq b_2$, las funciones $1$, $\sech^2(b_1 x)$ y $\sech^2(b_2 x)$ 
son linealmente independientes, por lo que tendríamos tres condiciones por 
ecuación. En particular, la primera ecuación exigiría $\alpha A_2^2 = 0$ y la 
segunda exigiría $\alpha A_1^2 = 0$, lo que fuerza $A_1 = 0$ o $A_2 = 0$ 
(solución trivial) siempre que $\alpha \neq 0$.

Por lo tanto, para obtener soluciones no triviales con acoplamiento $\alpha \neq 0$, 
debemos imponer
\begin{equation*}
    b_1 = b_2 =: b.
\end{equation*}
Bajo esta hipótesis, las ecuaciones se reducen a
\begin{align*}
    (b^2 + 1+\lambda_1) - (2b^2 + A_1^2)\sech^2(bx) - \alpha A_2^2\sech^2(bx) &= 0 \\
    (b^2 + 1+\lambda_2) - \alpha A_1^2\sech^2(bx) - (2b^2 + \beta A_2^2)\sech^2(bx) &= 0
\end{align*}
y ahora solo aparecen dos funciones linealmente independientes: $1$ y 
$\sech^2(bx)$. Identificando coeficientes obtenemos el sistema algebraico
\begin{align*}
    b^2 + 1 + \lambda_1 &= 0 \\
    2b^2 + A_1^2 + \alpha A_2^2 &= 0 \\
    b^2 + 1 + \lambda_2 &= 0 \\
    \alpha A_1^2 + 2b^2 + \beta A_2^2 &= 0
\end{align*}

Notamos que $\lambda_1=\lambda_2$, luego tenemos que 
\begin{align*}
    2b^2 =- A_1^2 - \alpha A_2^2  \\
    2b^2 =- \alpha A_1^2  -\beta A_2^2 
\end{align*}
Es decir, 

\begin{equation}
    A_1^2+\alpha A_2^2=\alpha A_1^2+\beta A_2^2 \\
    \implies A_1^2(1-\alpha)=A_2^2(\beta-\alpha) \\
    \implies A_1^2=A_2^2\frac{(\beta-\alpha)}{(1-\alpha)}
\end{equation}

Sustituimos en $A_1^2 + \alpha A_2^2 = -2b^2$:
\begin{equation*}
    \frac{\beta-\alpha}{1-\alpha}\,A_2^2 + \alpha A_2^2 = -2b^2
    \implies
    A_2^2\,\frac{\beta - \alpha + \alpha(1-\alpha)}{1-\alpha} = -2b^2
    \implies
    A_2^2\,\frac{\beta - \alpha^2}{1 - \alpha} = -2b^2,
\end{equation*}
de donde
\begin{equation}\label{eq:s4A2}
    A_2^2 = \frac{2b^2(\alpha - 1)}{\beta - \alpha^2}.
\end{equation}
Sustituyendo de vuelta:
\begin{equation}\label{eq:s4A1}
    A_1^2 = \frac{\beta-\alpha}{1-\alpha} \cdot \frac{2b^2(\alpha-1)}{\beta-\alpha^2}
    = \frac{2b^2(\alpha - \beta)}{\beta - \alpha^2}.
\end{equation}
La solución queda entonces
\begin{align*}
    \psi(t,x) &= A_1\sech(b(x-ct))\,
                 e^{i\left[(\lambda+\frac{c^2}{4})t - \frac{c}{2}x\right]}, \\
    \phi(t,x) &= A_2\sech(b(x-ct))\,
                 e^{i\left[(\lambda+\frac{c^2}{4})t - \frac{c}{2}x\right]},
\end{align*}
con parámetros
\begin{equation*}
    A_1^2 = \frac{2b^2(\alpha-\beta)}{\beta-\alpha^2}, \qquad
    A_2^2 = \frac{2b^2(\alpha-1)}{\beta-\alpha^2}, \qquad
    \lambda = -1 - b^2,
\end{equation*}
y $b > 0$ parámetro libre.

\subsection*{Condiciones de existencia}

Para que la solución sea no trivial se requiere $A_1^2 > 0$ y $A_2^2 > 0$. 
Como $b^2 > 0$, las condiciones se reducen a
\begin{equation*}
    \frac{\alpha - \beta}{\beta - \alpha^2} > 0 
    \qquad \text{y} \qquad 
    \frac{\alpha - 1}{\beta - \alpha^2} > 0.
\end{equation*}
Analizando por casos según el signo del denominador $\beta - \alpha^2$:

\begin{itemize}
    \item \textbf{Si $\beta - \alpha^2 > 0$:} se necesita $\alpha > \beta$ y 
    $\alpha > 1$. Pero $\alpha > 1 \implies \alpha^2 > \alpha$, y $\beta > \alpha^2 > \alpha$, 
    contradiciendo $\alpha > \beta$. \textbf{Imposible.}
    
    \item \textbf{Si $\beta - \alpha^2 < 0$:} se necesita $\alpha < \beta$ y 
    $\alpha < 1$. La condición $\alpha < \beta < \alpha^2$ con $\alpha < 1$ 
    tiene soluciones no vacías únicamente cuando $\alpha < \alpha^2$, es decir 
    $\alpha(\alpha-1) > 0$, lo que con $\alpha < 1$ fuerza $\alpha < 0$.
\end{itemize}
La región de existencia es por tanto
\begin{equation*}
    \alpha < 0 \qquad \text{y} \qquad \alpha < \beta < \alpha^2.
\end{equation*}


\begin{table}[H]
\centering
\renewcommand{\arraystretch}{2.5}
\resizebox{\textwidth}{!}{%
\begin{tabular}{|p{3cm}|c|c|c|c|}
\hline
\textbf{Condición} & $A_1^2$ & $A_2^2$ & $\lambda$ & \textbf{Solución} \\
\hline
$\alpha < 0$
\newline
\hspace*{-\tabcolsep}\rule{\dimexpr\linewidth+2\tabcolsep\relax}{0.4pt}
\newline $\alpha < \beta < \alpha^2$
& $\dfrac{2b^2(\alpha-\beta)}{\beta-\alpha^2}>0$
& $\dfrac{2b^2(\alpha-1)}{\beta-\alpha^2}>0$
& $-1-b^2$
& Ambas componentes \\
\hline
$\beta \geq \alpha^2$
& \multicolumn{2}{c|}{$A_1^2,A_2^2$ no pueden ser ambos positivos$^{*}$}
& ---
& No existe \\
\hline
$\beta < \alpha^2$
\newline
\hspace*{-\tabcolsep}\rule{\dimexpr\linewidth+2\tabcolsep\relax}{0.4pt}
\newline $\alpha \geq 0$
& \multicolumn{2}{c|}{$A_1^2,A_2^2$ no pueden ser ambos positivos$^{*}$}
& ---
& No existe \\
\hline
$\beta = \alpha$
\newline
\hspace*{-\tabcolsep}\rule{\dimexpr\linewidth+2\tabcolsep\relax}{0.4pt}
\newline $\alpha < 0$
& $0$
& $\dfrac{-2b^2}{\alpha}>0$
& $-1-b^2$
& Solo componente $\phi$ ($\psi\equiv0$) \\
\hline
$\beta = \alpha$
\newline
\hspace*{-\tabcolsep}\rule{\dimexpr\linewidth+2\tabcolsep\relax}{0.4pt}
\newline $\alpha > 0$, $\alpha\neq1$
& $0$
& $\dfrac{-2b^2}{\alpha}<0$
& ---
& No existe \\
\hline
$\alpha = 1$
\newline
\hspace*{-\tabcolsep}\rule{\dimexpr\linewidth+2\tabcolsep\relax}{0.4pt}
\newline $\beta \neq 1$
& $-2b^2<0$
& $0$
& ---
& No existe \\
\hline
\end{tabular}}
\caption{Casos de existencia de la Solución~4
(ansatz $u = A_1\operatorname{sech}(bx)$, $v = A_2\operatorname{sech}(bx)$, $b_1=b_2=b>0$ libre).
$^{*}$Si $\beta\geq\alpha^2$, el denominador $\beta-\alpha^2$ es positivo o nulo; suponiendo $A_1^2>0$ y $A_2^2>0$ se requeriría $\alpha>\beta$ y $\alpha>1$, pero entonces $\alpha^2>\alpha>\beta\geq\alpha^2$, contradicción. Si $\beta<\alpha^2$ y $\alpha\geq0$, $A_1^2>0$ exige $\alpha<\beta$ y $A_2^2>0$ exige $\alpha<1$; junto con $\alpha<\beta<\alpha^2$ esto requiere $\alpha<\alpha^2$, es decir $\alpha(\alpha-1)>0$, imposible si $0\leq\alpha<1$ y descartado si $\alpha\geq1$ por la primera condición. En ambos casos no existe solución no trivial con ambas componentes positivas.}
\label{tab:sol4_existencia}
\end{table}


Así la solución explicita es 



\begin{equation*}
    \psi(t,x) = \sqrt{\frac{2b^2(\beta-\alpha)}{\beta-\alpha^2}}\,
    \operatorname{sech}(b(x-ct))\,e^{i[(-(1+b^2)+\frac{c^2}{4})t-c\frac{x}{2}]}, \quad
    \phi(t,x) = \sqrt{\frac{2b^2(1-\alpha)}{\beta-\alpha^2}}\,
    \sech(b(x-ct))\,e^{i[(-(1+b^2)+\frac{c^2}{4})t-c\frac{x}{2}]},
\end{equation*}


Donde el grafico que representa la region de existencia dela solucion es 

\begin{figure}[H]
    \centering
    \includegraphics[width=0.5\textwidth]{TESIS/img/solucion4.png}
    \caption{Regiones de existencia de la Solución 1 (ansatz sech--sech).}
    \label{fig:regiones_sol1}
\end{figure}


\section{Energía y Momentum}

Usando la Energia encontrada en la Sección 1 del documento, tenemos que 
\begin{equation}\label{eq:energia}
    E(\psi,\phi) = \int_{\mathbb{R}} \frac{1}{2}|\psi_x|^2 + \frac{1}{2}|\phi_x|^2
    + \frac{1}{4}(1 - |\psi|^2 - |\phi|^2)^2
    + \frac{\alpha - 1}{2}|\psi|^2|\phi|^2
    + \frac{\beta - 1}{4}|\phi|^4 \, dx,
\end{equation}

Ademas el momentum esta determinado por


\begin{equation*}
    P=\frac{1}{2}\int_\R Re(i\psi \overline{\psi_x})dx+\frac{1}{2}\int_\R Re(i\phi \overline{\phi_x})dx
\end{equation*}



\subsection{Energia Solucion 1.2}
Donde nos interesa calcular la energía de la Solución 1.2

\begin{align*}
     u = A_1\tanh(bx) + iB_1 \\
     v = A_2\sech(bx)
\end{align*}

Donde tenemos que  
\begin{align*}
\psi(t,x) &=
\sqrt{\frac{2(\beta-\alpha^2)-c^2(\beta-\alpha)}{2(\beta-\alpha^2)}}
\tanh\!\left(\sqrt{\frac{2(\beta-\alpha^2)-c^2(\beta-\alpha)}{4(\beta-\alpha)}}\,(x-ct)\right)
+\,i\sqrt{\frac{c^2(\beta-\alpha)}{2(\beta-\alpha^2)}}
\\[12pt]
\phi(t,x) &=
\sqrt{\frac{(\alpha-1)\bigl[2(\beta-\alpha^2)-c^2(\beta-\alpha)\bigr]}{2(\beta-\alpha)(\beta-\alpha^2)}}
\;\operatorname{sech}\!\left(\sqrt{\frac{2(\beta-\alpha^2)-c^2(\beta-\alpha)}{4(\beta-\alpha)}}\,(x-ct)\right)
\\[6pt]
&\quad\times
\exp\!\left(i\left[\left(\alpha-1-\frac{2(\beta-\alpha^2)-c^2(\beta-\alpha)}{4(\beta-\alpha)}+\frac{c^2}{4}\right)t
-\frac{c}{2}\,x\right]\right)
\end{align*}
Debemos remplazar esto en la ecuacion de energia. Denotamos como $\mathcal{N}=2(\beta-\alpha^2)-c^2(\beta-\alpha)$ y $$ \kappa= \sqrt{\frac{2(\beta-\alpha^2)-c^2(\beta-\alpha)}{4(\beta-\alpha)}}$$ para simplificación de cálculos

\begin{align*}
    |\psi|^2&=\frac{\mathcal{N}}{2(\beta-\alpha^2)}
\tanh^2\!\left(\kappa\,(x-ct)\right)+\frac{c^2(\beta-\alpha)}{2(\beta-\alpha^2)} \\
            &=\frac{\mathcal{N}-\mathcal{N}\sech^2(\kappa(x-ct))+c^2(\beta-\alpha)}{2(\beta-\alpha^2)}\\
            &=\frac{2(\beta-\alpha^2)-\mathcal{N}\sech^2(\kappa(x-ct))}{2(\beta-\alpha^2)} \\
            &=1-\frac{\mathcal{N}\sech^2(\kappa(x-ct))}{2(\beta-\alpha^2)}
\end{align*}

Ahora desarrollamos $|\phi|$

\begin{align*}
    |\phi|^2=\frac{(\alpha-1)\mathcal{N}}{2(\beta-\alpha)(\beta-\alpha^2)}\sech^2(\kappa(x-ct))
\end{align*}

\begin{align*}
    |\phi|^4=\frac{(\alpha-1)^2\mathcal{N}^2}{4(\beta-\alpha)^2(\beta-\alpha^2)^2}\sech^4(\kappa(x-ct))
\end{align*}

Calculamos el resto de expresiones
\begin{equation*}
    1-|\phi|^2-|\psi|^2=\frac{\mathcal{N}(\beta-2\alpha+1)}{2(\beta-\alpha^2)(\beta-\alpha)}\sech^2(\kappa(x-ct))
\end{equation*}

\begin{equation*}
    |\phi|^2|\psi|^2=\frac{(\alpha-1)\mathcal{N}[2(\beta-\alpha^2)-\mathcal{N}\sech^2(\kappa(x-ct)]}{2(\beta-\alpha^2)^2(\beta-\alpha)}\sech^2(\kappa(x-ct))
\end{equation*}

\begin{equation*}
    |\psi_x|^2=\frac{\mathcal{N}^2}{8(\beta-\alpha^2)(\beta-\alpha)}\sech^4(\kappa (x-ct))
\end{equation*}

\begin{equation*}
    |\phi_x|^2=\frac{(\alpha-1)\mathcal{N}}{2(\beta-\alpha)(\beta-\alpha^2)}\sech^2(\kappa(x-ct))\left[\kappa^2\tanh^2(\kappa(x-ct))+\frac{c^2}4{}\right] 
\end{equation*}

Ahora desarrollaremos las integrales

\begin{align*}
    \int_{\mathbb{R}} \frac{1}{2}|\psi_x|^2&= \int_{\mathbb{R}}  \frac{\mathcal{N}^2}{16(\beta-\alpha^2)(\beta-\alpha)}\sech^4(\kappa (x-ct)) =\frac{\mathcal{N}^2}{16(\beta-\alpha^2)(\beta-\alpha)}\cdot\frac{4}{3}\cdot \sqrt{\frac{4(\beta-\alpha)}{2(\beta-\alpha^2)-c^2(\beta-\alpha)}} \\
    &=\frac{\mathcal{N}^\frac{3}{2}}{6(\beta-\alpha^2)(\beta-\alpha)^{1/2}}
\end{align*}

\begin{align*}
\frac{1}{2}\int_{\mathbb{R}}|\phi_x|^2\,dx &= \frac{1}{2}\frac{(\alpha-1)\mathcal{N}}{2(\beta-\alpha)(\beta-\alpha^2)}\left[\left(\kappa^2+\frac{c^2}{4}\right)\int_{\mathbb{R}}\operatorname{sech}^2(\kappa\xi)\,d\xi - \kappa^2\int_{\mathbb{R}}\operatorname{sech}^4(\kappa\xi)\,d\xi\right]\\[8pt]
&= \frac{(\alpha-1)\mathcal{N}}{4(\beta-\alpha)(\beta-\alpha^2)}\left[\left(\kappa^2+\frac{c^2}{4}\right)\frac{2}{\kappa} - \kappa^2\cdot\frac{4}{3\kappa}\right]\\[8pt]
&= \frac{(\alpha-1)\mathcal{N}}{4(\beta-\alpha)(\beta-\alpha^2)}\left[2\kappa+\frac{c^2}{2\kappa}-\frac{4\kappa}{3}\right]\\[8pt]
&= \frac{(\alpha-1)\mathcal{N}}{4(\beta-\alpha)(\beta-\alpha^2)}\left[\frac{2\kappa}{3}+\frac{c^2}{2\kappa}\right]\\[8pt]
&= \frac{(\alpha-1)\mathcal{N}}{4(\beta-\alpha)(\beta-\alpha^2)}\left[\frac{\sqrt{\mathcal{N}}}{3\sqrt{\beta-\alpha}}+\frac{c^2\sqrt{\beta-\alpha}}{\sqrt{\mathcal{N}}}\right]
\end{align*}

\begin{align*}
\frac{\alpha-1}{2}\int_{\mathbb{R}}|\phi|^2|\psi|^2\,dx &= \frac{\alpha-1}{2}\cdot\frac{(\alpha-1)\mathcal{N}}{2(\beta-\alpha^2)^2(\beta-\alpha)}\left[2(\beta-\alpha^2)\int_{\mathbb{R}}\operatorname{sech}^2(\kappa\xi)\,d\xi - \mathcal{N}\int_{\mathbb{R}}\operatorname{sech}^4(\kappa\xi)\,d\xi\right]\\[8pt]
&= \frac{(\alpha-1)^2\mathcal{N}}{4(\beta-\alpha^2)^2(\beta-\alpha)}\left[2(\beta-\alpha^2)\cdot\frac{2}{\kappa} - \mathcal{N}\cdot\frac{4}{3\kappa}\right]\\[8pt]
&= \frac{(\alpha-1)^2\mathcal{N}}{4(\beta-\alpha^2)^2(\beta-\alpha)}\cdot\frac{4}{\kappa}\left[(\beta-\alpha^2) - \frac{\mathcal{N}}{3}\right]\\[8pt]
&= \frac{2(\alpha-1)^2\sqrt{\mathcal{N}}}{(\beta-\alpha^2)^2\sqrt{\beta-\alpha}}\left[(\beta-\alpha^2)-\frac{\mathcal{N}}{3}\right]
\end{align*}

\begin{align*}
\frac{\beta-1}{4}\int_{\mathbb{R}}|\phi|^4\,dx &= \frac{\beta-1}{4}\cdot\frac{(\alpha-1)^2\mathcal{N}^2}{4(\beta-\alpha)^2(\beta-\alpha^2)^2}\int_{\mathbb{R}}\operatorname{sech}^4(\kappa\xi)\,d\xi\\[8pt]
&= \frac{(\beta-1)(\alpha-1)^2\mathcal{N}^2}{16(\beta-\alpha)^2(\beta-\alpha^2)^2}\cdot\frac{4}{3\kappa}\\[8pt]
&= \frac{(\beta-1)(\alpha-1)^2\mathcal{N}^2}{12(\beta-\alpha)^2(\beta-\alpha^2)^2\kappa}\\[8pt]
&= \frac{(\beta-1)(\alpha-1)^2\mathcal{N}^{3/2}}{6(\beta-\alpha)^{3/2}(\beta-\alpha^2)^2}
\end{align*}

\begin{align*}
\frac{1}{4}\int_{\mathbb{R}}\left(1-|\phi|^2-|\psi|^2\right)^2\,dx &= \frac{1}{4}\cdot\frac{\mathcal{N}^2(\beta-2\alpha+1)^2}{4(\beta-\alpha^2)^2(\beta-\alpha)^2}\int_{\mathbb{R}}\operatorname{sech}^4(\kappa\xi)\,d\xi\\[8pt]
&= \frac{\mathcal{N}^2(\beta-2\alpha+1)^2}{16(\beta-\alpha^2)^2(\beta-\alpha)^2}\cdot\frac{4}{3\kappa}\\[8pt]
&= \frac{\mathcal{N}^2(\beta-2\alpha+1)^2}{12(\beta-\alpha^2)^2(\beta-\alpha)^2\kappa}\\[8pt]
&= \frac{\mathcal{N}^{3/2}(\beta-2\alpha+1)^2}{6(\beta-\alpha^2)^2(\beta-\alpha)^{3/2}}
\end{align*}

Usando lo anterior podemos escribir la Energia de la siguiente manera
\begin{align*}
   E(\psi,\phi)&= \frac{\mathcal{N}^\frac{3}{2}}{6(\beta-\alpha^2)(\beta-\alpha)^{1/2}}+\frac{(\alpha-1)\mathcal{N}}{4(\beta-\alpha)(\beta-\alpha^2)}\left[\frac{\sqrt{\mathcal{N}}}{3\sqrt{\beta-\alpha}}+\frac{c^2\sqrt{\beta-\alpha}}{\sqrt{\mathcal{N}}}\right]+\frac{2(\alpha-1)^2\sqrt{\mathcal{N}}}{(\beta-\alpha^2)^2\sqrt{\beta-\alpha}}\left[(\beta-\alpha^2)-\frac{\mathcal{N}}{3}\right]\\[8pt]
   &+\frac{(\beta-1)(\alpha-1)^2\mathcal{N}^{3/2}}{6(\beta-\alpha)^{3/2}(\beta-\alpha^2)^2}+\frac{\mathcal{N}^{3/2}(\beta-2\alpha+1)^2}{6(\beta-\alpha^2)^2(\beta-\alpha)^{3/2}}
\end{align*}

Desarrollamos esta expresión 

\begin{align*}
E(\psi,\phi) &= \frac{\sqrt{2(\beta-\alpha^2)-c^2(\beta-\alpha)}}{(\beta-\alpha^2)\sqrt{\beta-\alpha}}\left\{(\alpha-1)\left[\frac{c^2}{4}+2(\alpha-1)\right]\right.\\[8pt]
&\quad+\left.[2(\beta-\alpha^2)-c^2(\beta-\alpha)]\left[\frac{1}{6}+\frac{\alpha-1}{12(\beta-\alpha)}+\frac{4\alpha^3-(3\beta+5)\alpha^2+2(\beta+1)\alpha+\beta(\beta-1)}{6(\beta-\alpha)(\beta-\alpha^2)}\right]\right\}
\end{align*}

Recordamos que $\mathcal{N} = 2(\beta-\alpha^2)-c^2(\beta-\alpha)$ y $\kappa = \sqrt{\dfrac{2(\beta-\alpha^2)-c^2(\beta-\alpha)}{4(\beta-\alpha)}}$. Para que la expresión de energia esté bien definida y la solución exista, se requieren las siguientes condiciones sobre los parámetros:

\begin{itemize}
    \item $\beta - \alpha > 0$: garantiza que $\kappa$ sea real, pues $(\beta-\alpha)$ aparece en el denominador de $\kappa^2$. Si $\beta - \alpha < 0$, entonces aunque el numerador $\mathcal{N}$ fuera positivo, el cociente $\dfrac{\mathcal{N}}{4(\beta-\alpha)}$ sería negativo, haciendo que $\kappa$ sea imaginario y por tanto la solución deja de ser una función real localizada. Si $\beta - \alpha = 0$, el denominador se anula y $\kappa$ queda indefinido.
    
    \item $\beta - \alpha^2 > 0$: garantiza que $\mathcal{N} > 0$ para al menos un rango de valores de $c$. En efecto, bajo la condición $\beta - \alpha > 0$, la condición $\mathcal{N} > 0$ equivale a
    \begin{equation*}
        c^2 < \frac{2(\beta-\alpha^2)}{\beta-\alpha},
    \end{equation*}
    por lo que el parámetro $c$ queda restringido al intervalo $c \in (-c_{\max}, c_{\max})$, donde
    \begin{equation*}
        c_{\max} = \sqrt{\frac{2(\beta-\alpha^2)}{\beta-\alpha}}.
    \end{equation*}
    Si $\beta - \alpha^2 \leq 0$, entonces $\dfrac{2(\beta-\alpha^2)}{\beta-\alpha} \leq 0$ y no existe ningún valor real de $c$ que satisfaga dicha desigualdad, de modo que $\mathcal{N} \leq 0$ para todo $c \in \mathbb{R}$ y la solución deja de existir.
\end{itemize}

\subsection{Energía Solución 4}

Denotamos $ \xi =x-ct$ luego tenemos que 

\begin{align*}
    |\psi|^2 &=A_1^2\sech^2(b\xi) \\
    |\phi|^2 &=A_2^2\sech^2(b\xi) \\
    |\psi_x|^2 &= A_1^2\sech^2(b\xi)\left [b^2\tanh^2(b\xi)+\frac{c^2}{4} \right] \\
    |\psi_x|^2 &= A_2^2\sech^2(b\xi)\left [b^2\tanh^2(b\xi)+\frac{c^2}{4} \right]
\end{align*}

Recordamos las siguientes integrales 

\begin{equation*}
  \int_\R \sech^2(b\xi)d\xi =\frac{2}{b} \qquad \int_\R \sech^2(b\xi)\tanh ^2(b\xi)d\xi =\frac{2}{3b} \qquad \int_\R \sech^4(b\xi)d\xi =\frac{4}{3b}
\end{equation*}

Luego asi tenemos que los resultados de los terminos de energia son

\begin{align*}
    &\int_{\mathbb{R}} \frac{1}{2}|\psi_x|^2 + \frac{1}{2}|\phi_x|^2 dx= \frac{1}{2}[A_1^2+A_2^2]\int_\R [b^2\sech^2(b\xi)\tanh^2(b\xi)+\frac{c^2} {4}\sech^2(b\xi)]dx =\frac{A_1^2+A_2^2}{2}\left [\frac{2b}{3} + \frac{c^2}{2b}\right ] \\
    &\int_\R \frac{1}{4}(1 - |\psi|^2 - |\phi|^2)^2 dx=[A_1^2+A_2^2]\left [-\frac{1}{b}+\frac{A_1^2+A_2^2d}{3b}\right] \\
     &\int_\R \frac{\alpha - 1}{2}|\psi|^2|\phi|^2 dx =\frac{2(\alpha - 1)}{3b}A_1^2A_2^2 \\
     &\int_\R \frac{\beta - 1}{4}|\phi|^4 \, dx =\frac{\beta - 1}{3b}A_2^4
\end{align*}

Ahora sumando todo lo anterior, nos definimos $S:=A_1^2+A_2^2$ obtenemos que 

\begin{align*}
   E(\psi,\phi)&= \frac{Sb}{3}+\frac{Sc^2}{4b}-\frac{S}{b}+\frac{S^2}{3b}+\frac{1}{3b}[2(\alpha-1)A_1^2A_2^2+(\beta-1)A_2^4] \\
   E(\psi,\phi)&= S\left [ \frac{b}{3} + \frac{c^2-4}{4b}\right]+\frac{S^2}{3b} +\frac{1}{3b}[2(\alpha-1)A_1^2A_2^2+(\beta-1)A_2^4] 
\end{align*}
Donde tenemos que $S^2=A_1^4+2A_1^2A_2^2+A_2^4$
\begin{align*}
   E(\psi,\phi)&= S\left [ \frac{b}{3} + \frac{c^2-4}{4b}\right]+\frac{A_1^4+2A_1^2A_2^2+A_2^4}{3b} +\frac{1}{3b}[2(\alpha-1)A_1^2A_2^2+(\beta-1)A_2^4] \\
    E(\psi,\phi)&= S\left [ \frac{b}{3} + \frac{c^2-4}{4b}\right] +\frac{1}{3b}[A_1^4+2\alpha A_1^2A_2^2+\beta A_2^4] 
\end{align*}

Usando los valores encontrados de $A_1^2$ y $A_2^2$ notamos que $S=\frac{2b^2(2\alpha-\beta-1)}{\beta-\alpha^2}$. Calculamos $A_1^4 + 2\alpha A_1^2A_2^2 + \beta A_2^4$ sustituyendo nuestros valores 

\begin{align*}
A_1^4 + 2\alpha A_1^2A_2^2 + \beta A_2^4 
&= \frac{4b^4}{(\beta-\alpha^2)^2}\left[(\alpha-\beta)^2 + 2\alpha(\alpha-\beta)(\alpha-1) 
   + \beta(\alpha-1)^2\right]
\end{align*}
Expandimos el corchete término a término:
\begin{align*}
(\alpha-\beta)^2 &= \alpha^2 - 2\alpha\beta + \beta^2 \\
2\alpha(\alpha-\beta)(\alpha-1) &= 2\alpha^3 - 2\alpha^2 - 2\alpha^2\beta + 2\alpha\beta \\
\beta(\alpha-1)^2 &= \alpha^2\beta - 2\alpha\beta + \beta
\end{align*}
Sumando:
\begin{equation*}
(\alpha-\beta)^2 + 2\alpha(\alpha-\beta)(\alpha-1) + \beta(\alpha-1)^2 
= 2\alpha^3 - \alpha^2 - \alpha^2\beta - 2\alpha\beta + \beta^2 + \beta
\end{equation*}
que se factoriza como
\begin{equation*}
= (\alpha^2-\beta)(2\alpha-\beta-1) = -(\beta-\alpha^2)(2\alpha-\beta-1).
\end{equation*}
Por lo tanto
\begin{equation*}
A_1^4 + 2\alpha A_1^2A_2^2 + \beta A_2^4 
= \frac{4b^4}{(\beta-\alpha^2)^2}\cdot(-(\beta-\alpha^2))(2\alpha-\beta-1) 
= \frac{-4b^4(2\alpha-\beta-1)}{(\beta-\alpha^2)}
\end{equation*}
Ahora sustituimos esto en la energia y obtenemos que 
\begin{equation}
    E(\psi,\phi)=\frac{2b^2(2\alpha-\beta-1)}{\beta-\alpha^2}\left [ \frac{b}{3} + \frac{c^2-4}{4b}\right]+\frac{1}{3b}\frac{-4b^4(2\alpha-\beta-1)}{(\beta-\alpha^2)}
\end{equation}

Factorizamos esta expresion y obtenemos que 
\begin{equation*}
    E(\psi,\phi)=\frac{b(2\alpha-\beta-1)}{\beta-\alpha^2}\left [\frac{2b^2}{3} +\frac{c^2-4}{2}-\frac{4b^2}{3}\right]
\end{equation*}

Asi finalmente obtenemos que la energia es 
\begin{equation*}
    E(\psi,\phi)=\frac{b(2\alpha-\beta-1)(3c^2-4b^2-12)}{6(\beta-\alpha^2)}
\end{equation*}

\subsection{Momentum Solución 4}

Usamos las notaciones establecidas en el calculo de la energía y que $\theta=e^{i\left[(\lambda+\frac{c^2}{4})t - \frac{c}{2}x\right]}$ donde 



\begin{equation*}
    \psi_x=A_1e^{i\theta}\sech(b\xi)\left [-b\tanh(b\xi)-i\frac{c}{2} \right]
\end{equation*}

\begin{equation*}
   \overline{ \psi_x}=A_1e^{i\theta}\sech(b\xi)\left [-b\tanh(b\xi)+i\frac{c}{2} \right]
\end{equation*}

Procedemos con los cálculos 
\begin{align*}
   \psi \overline{ \psi_x}&=A_1\sech(b\xi)e^{i\theta}A_1e^{i\theta}\sech(b\xi)\left [-b\tanh(b\xi)+i\frac{c}{2} \right] \\
   &= A_1^2\sech^2(b\xi)\left [-b\tanh(b\xi)+i\frac{c}{2} \right] 
\end{align*}

Asi tenemos que $Re(i\psi \overline{ \psi_x})=-A_1^2\sech^2(b\xi)\cdot \frac{c}{2}$. Como los cálculos son análogos para $\phi$, se tiene que  $$Re(i\phi \overline{ \phi_x})=-A_2^2\sech^2(b\xi)\cdot \frac{c}{2}$$ remplazando esto en la ecuación de Potencia

\begin{align*}
    P(\phi,\psi)&=\frac{1}{2}\int_\R Re(i\psi \overline{\psi_x})dx+\frac{1}{2}\int_\R Re(i\phi \overline{\phi_x})dx \\\hspace{7mm}
    &= \frac{-A_1^2c}{4}\int_\R \sech^2(b\xi)dx +\frac{-A_2^2c}{4}\int_\R \sech^2(b\xi)dx 
\end{align*}

Finalmente obtenemos que

\begin{equation}
    P(\phi,\psi)=\frac{-cb(2\alpha-\beta-1)}{\beta-\alpha^2}
\end{equation}

Con el resultado encontrado notamos que el momentum se puede trabajar como una funcion de c y fijando b.


\section{Sponge Factor}
En las secciones anteriores hemos construido soluciones exactas del tipo bright-brigh, dark-dark, black soliton que satisfacen nuestro sistema bajo condiciones ideales, perfil exacto, parametros exactos, dominio infinito. Pero en la practica, nos surgen preguntas que no se pueden responder con la solución exacta.

\hspace{2mm}

¿Nuestro soliton es estable?, ¿Si perturbamos la condición inicial, nuestro soliton se mantiene cerca del soliton original en el tiempo o este cambia totalmente?

\hspace{2mm}

Para responder estas preguntas necesitamos entender como evoluciona numéricamente el sistema bajo condiciones iniciales perturbadas, el reto ahora es elegir un esquema numérico que represente fielmente las propiedades de nuestro sistema y los solitones.

\hspace{2mm}


\subsection{Conservación de la norma $L^2$}
Una de las propiedades fundamentales para determinar si nuestra simulación numérica representa las propiedades de los solitones consiste en la conservación de la norma $L^2$, que físicamente es la conservación de la densidad. Definimos la densidad de $\psi$ es análogo para $\phi$

\begin{equation}
    N_\psi(t)=\int_\R|\psi(t,x)|^2 dx
\end{equation}


Entonces nuestra simulación debe cumplir la condición
\begin{equation}
\frac{dN_\psi}{dt}=0
\end{equation}

Es decir, las soluciones del sistema satisfacen $N_\psi(t) = N_\psi(0)$ y 
$N_\phi(t) = N_\phi(0)$ para todo $t \geq 0$.

\hspace{2mm}

\textbf{Demostracion}


Por el teorema de Leibniz para derivación bajo el signo integral, podemos 
intercambiar el orden de la derivada y la integral:

\begin{equation*}
    \frac{dN_\psi}{dt} = \int_{\mathbb{R}}\partial_t|\psi(t,x)|^2\,dx.
\end{equation*}

Como $|\psi|^2 = \psi\cdot\bar\psi$, aplicamos la regla del producto

\begin{equation*}
    \partial_t|\psi|^2 = (\partial_t\psi)\bar\psi + \psi(\partial_t\bar\psi).
\end{equation*}

De la primera ecuación del sistema, despejamos $\partial_t\psi$ usando 
$1/i = -i$:

\begin{equation*}
    \partial_t\psi = -i\left[\partial_{xx}\psi + 
    \psi(1-|\psi|^2-\alpha|\phi|^2)\right].
\end{equation*}

Tomando el conjugado complejo, y usando que la no linealidad 
$1-|\psi|^2-\alpha|\phi|^2$ es real:

\begin{equation*}
    \partial_t\bar\psi = i\left[\partial_{xx}\bar\psi + 
    \bar\psi(1-|\psi|^2-\alpha|\phi|^2)\right].
\end{equation*}

Sustituyendo en la regla del producto y expandiendo:

\begin{equation*}
    \partial_t|\psi|^2 = -i\bar\psi\,\partial_{xx}\psi 
    -i|\psi|^2(1-|\psi|^2-\alpha|\phi|^2) 
    + i\psi\,\partial_{xx}\bar\psi 
    +i|\psi|^2(1-|\psi|^2-\alpha|\phi|^2)
\end{equation*}


Se cancelan terminos, donde esta cancelación es válida para cualquier no linealidad real 
$f(|\psi|^2,|\phi|^2)$, independientemente de los valores de 
$\alpha$ y $\beta$. Asi obtenemos:

\begin{equation*}
    \partial_t|\psi|^2 = i\left(\psi\,\partial_{xx}\bar\psi 
    - \bar\psi\,\partial_{xx}\psi\right).
\end{equation*}

Integrando en $x \in \mathbb{R}$ y aplicando integración por partes 
con $u = \bar\psi$, $v' = \partial_x\psi$:

\begin{equation*}
    \int_{\mathbb{R}}\bar\psi\,\partial_{xx}\psi\,dx = 
    \underbrace{[\bar\psi\,\partial_x\psi]_{-\infty}^{+\infty}}_{=\,0} 
    - \int_{\mathbb{R}}|\partial_x\psi|^2\,dx 
    = -\int_{\mathbb{R}}|\partial_x\psi|^2\,dx,
\end{equation*}

donde el término de borde es cero porque para el bright soliton 
$\psi \to 0$ en el infinito, y para el black/dark soliton 
$\partial_x\psi \to 0$ en el infinito. El mismo argumento da

$$\int_{\mathbb{R}}\psi\,\partial_{xx}\bar\psi\,dx = 
-\int_{\mathbb{R}}|\partial_x\psi|^2\,dx$$

Entonces uniendo lo anterior obtenemos que

\begin{equation*}
    \frac{dN_\psi}{dt} = i\left(
    -\int_{\mathbb{R}}|\partial_x\psi|^2\,dx + 
    \int_{\mathbb{R}}|\partial_x\psi|^2\,dx\right) = 0,
\end{equation*}

\hspace{2mm}


de donde $N_\psi(t) = N_\psi(0)$ para todo $t \geq 0$. El mismo 
argumento se aplica a $\phi$.

\hspace{2mm}

--------------------------------

\subsection{Discretización}

Discretizamos el espacio como $x_j = jh$ y el tiempo como 
$t_n = n\,\delta t$, y denotamos $\psi_j^n \approx \psi(t_n,x_j)$ 
y $\phi_j^n \approx \phi(t_n,x_j)$. Asi obtenemos que 

\begin{equation*}
    \partial_t\psi(t_n,x_j)\approx \frac{\psi(t_{n+1},x_j)-\psi(t_{n},x_j)}{\delta t} = \frac{\psi_j^{n+1}-\psi_j^n}{\delta t}
\end{equation*}


Analicemos la aproximación del operador $\partial_xx$ usamos la expansion de taylor, tenemos 

\begin{align*}
    \psi(x_j+h)=\psi(x_j)+ h\partial_x \psi(x_j)+\frac{h^2}{2}\partial_{xx}\psi(x_j)+ \frac{h^3}{3!}\partial_{xxx}\psi(x_j)+..\\
     \psi(x_j-h)=\psi(x_j)-h\partial_x \psi(x_j)+\frac{h^2}{2}\partial_{xx}\psi(x_j)- \frac{h^3}{3!}\partial_{xxx}\psi(x_j)+..\\
\end{align*}

Sumamos estas ecuaciones y notamos que los términos con potencia impar se cancelan por lo que obtenemos

\begin{equation*}
    \psi(x_j+h)+ \psi(x_j-h)=2\psi(x_j)+h^2\partial_{xx}\psi(x_j)+ \mathcal{O}(h^4)
\end{equation*}

Despejamos $\partial_{xx}\psi(x_j)$ y obtenemos 
\begin{equation*}
    \partial_{xx}\psi(x_j)\approx\frac{\psi(x_j+h)+\psi(x_j-h)+2\psi(x_j)}{h^2}
\end{equation*}

Usando nuestra notacion tenemos que 

\begin{itemize}
    \item $\psi(x_j+h)=\psi_{j+1}^n$
    \item $\psi(x_j-h)=\psi_{j-1}^n$
    \item $\psi(x_j)=\psi_j^n$
\end{itemize}

remplazamos esto en lo anterior y obtenemos la siguiente ecuacion

\begin{equation*}
    \partial_{xx}\psi(t_n,x_j)\approx\frac{\psi_{j+1}^n-2\psi_j^n+\psi_{j-1}^n}{h^2}
\end{equation*}



Definimos el operador de diferencias finitas centradas:
\begin{equation*}
    (L_xu)_j:= \frac{u_{j+1}-2u_j+u_{j-1}}{h^2}\approx \partial_{xx}u(x_j)
\end{equation*}


\subsection{Sponge factor}
Al simular numericamente nuestra EDP en el dominio infinito $\R$, debemos truncar este dominio a un intervalo finito $[-L,L]$, pero esto nos genera un problema, debido a que se generan bordes artificiales que no están presentes en nuestro problema original. Esto provoca que nuestro solitón, al llegar a estas fronteras, rebote y vuelva al interior del dominio, lo que provocara que se contamine nuestra simulacion.


\hspace{2mm}

Una forma de atenuar o contrarestar este fenómeno, se conoce como \textbf{Sponge Factor} que es un termino que se establece cerca de los bordes con el fin de absorber la energia del soliton, asi simulando un dominio infinito.

\hspace{2mm}

Siguiendo a Di Menza y Gallo, reemplazamos $i\partial_t u$ por $(i-\eta(x))\partial_t u$, donde $\eta(x)$ es una función que se anula en el interior del dominio y crece cerca de los bordes, absorbiendo la energía 
antes de que llegue a la frontera. Así nuestro sistema queda como

\begin{align}
    (i-\eta(x) )\partial_t \psi&=\partial_{xx}\psi + \psi (1-|\psi|^2-\alpha|\phi|^2)  \\
         (i-\eta(x) ) \partial_t \phi&=\partial_{xx}\phi + \phi(1-\alpha|\phi|^2-\beta|\phi|^2) 
\end{align}

Escogemos $\eta(x)$ como

\begin{equation}
    \eta(x)=exp \left ( -\frac{(|x|-L)^2}{4}\right)
\end{equation}





Esta función satisface $\eta(x) \approx 0$ en el interior 
(donde $|x| \ll L$) y $\eta(\pm L) = 1$ en los bordes, con 
transición suave entre ambos valores. Donde $\eta \approx 0$ 
se recupera exactamente la ecuación original; donde $\eta > 0$ 
el operador $(i-\eta)\partial_t$ introduce disipación que 
amortigua la solución antes de que alcance el borde.

\subsection{Esquema de Crank-Nicolson}
Seguiremos el metodo de Crank-Nicolson, para aproximar el lado derecho de la ecuación en el tiempo intermedio $t_{n+1/2}$, tomando el 
promedio entre $t_n$ y $t_{n+1}$. Definimos:

\begin{equation*}
   F(u)= u^{n+1/2} := \frac{u^n + u^{n+1}}{2}.
\end{equation*}

El esquema completo con sponge factor es:

\begin{align}
    (i-\eta(x_j))\frac{\psi_j^{n+1}-\psi_j^n}{\delta t} + 
    (L_x\psi^{n+1/2})_j + 
    \psi_j^{n+1/2}(1-|\psi_j^{n+1/2}|^2-
    \alpha|\phi_j^{n+1/2}|^2) &= 0 \\
    (i-\eta(x_j))\frac{\phi_j^{n+1}-\phi_j^n}{\delta t} + 
    (L_x\phi^{n+1/2})_j + 
    \phi_j^{n+1/2}(1-\alpha|\psi_j^{n+1/2}|^2-
    \beta|\phi_j^{n+1/2}|^2) &= 0
\end{align}

donde $(L_x\psi^{n+1/2})_j = \dfrac{\psi_{j+1}^{n+1/2} - 
2\psi_j^{n+1/2} + \psi_{j-1}^{n+1/2}}{h^2}$ es la diferencia 
finita centrada aplicada al promedio temporal.

\hspace{2mm}

Para mayor claridad las ecuaciones anteriores sin usar funciones alternativas para la simplificacion de terminos, se ven representadas en su formato expandido de la siguiente manera


{\footnotesize
\begin{align*}
    (i-\eta(x_j))\frac{\psi_j^{n+1}-\psi_j^n}{\delta t} + 
    \frac{\dfrac{\psi_{j+1}^n+\psi_{j+1}^{n+1}}{2} - 
    2\dfrac{\psi_j^n+\psi_j^{n+1}}{2} + 
    \dfrac{\psi_{j-1}^n+\psi_{j-1}^{n+1}}{2}}{h^2} + 
    \frac{\psi_j^n+\psi_j^{n+1}}{2}\left(1-
    \left|\frac{\psi_j^n+\psi_j^{n+1}}{2}\right|^2-
    \alpha\left|\frac{\phi_j^n+\phi_j^{n+1}}{2}\right|^2\right) &= 0 \\
    (i-\eta(x_j))\frac{\phi_j^{n+1}-\phi_j^n}{\delta t} + 
    \frac{\dfrac{\phi_{j+1}^n+\phi_{j+1}^{n+1}}{2} - 
    2\dfrac{\phi_j^n+\phi_j^{n+1}}{2} + 
    \dfrac{\phi_{j-1}^n+\phi_{j-1}^{n+1}}{2}}{h^2} + 
    \frac{\phi_j^n+\phi_j^{n+1}}{2}\left(1-
    \alpha\left|\frac{\psi_j^n+\psi_j^{n+1}}{2}\right|^2-
    \beta\left|\frac{\phi_j^n+\phi_j^{n+1}}{2}\right|^2\right) &= 0
\end{align*}}

\hspace{2mm}

Como $u^{n+1}$ 
aparece en ambos lados a través de $u^{n+1/2}$, este es un 
sistema no lineal en $(\psi^{n+1},\phi^{n+1})$ que se resuelve 
en cada paso de tiempo mediante \textbf{iteración de punto fijo}: 
se inicializa $\psi^{(0)}=\psi^n$, $\phi^{(0)}=\phi^n$, y se 
itera hasta convergencia.

\subsection{Por qué Crank-Nicolson conserva la norma discreta}

Para verificar que el esquema conserva $N_\psi^n$, usamos el 
\textbf{análisis de Von Neumann}: estudiamos cómo evoluciona 
un solo modo de Fourier $\psi_j^n = g^n e^{ikjh}$ bajo el 
esquema lineal. Esto es válido porque el esquema puede 
escribirse como $\psi^{n+1} = M\psi^n$ donde $M$ es una 
\textbf{matriz circulante} (cada fila es un desplazamiento 
de la anterior, pues el esquema trata todos los puntos $x_j$ 
de la misma forma). Por el siguiente teorema, esto garantiza 
que los modos de Fourier son sus autovectores:

\begin{theorem}
Los autovectores de toda matriz circulante de tamaño $N$ son 
los modos de Fourier discretos $v_k = (e^{ikx_{-J}},\ldots,
e^{ikx_J})^\top$, con autovalores $g_k$ que dependen de $k$.
\end{theorem}

Como $\{v_k\}$ forma una base ortonormal, cualquier condición 
inicial se escribe $\psi^0 = \sum_k c_k v_k$, y por linealidad 
cada modo evoluciona independientemente:

\begin{equation*}
    \psi^n = \sum_k c_k g_k^n v_k \implies 
    N_\psi^n = h\sum_k|c_k|^2|g_k|^{2n}.
\end{equation*}

Para que $N_\psi^n = N_\psi^0$ para todo $n$, es necesario 
y suficiente que $|g_k|=1$ para todo $k$.

Sustituyendo el modo $\psi_j^n = g^n e^{ikjh}$ en el esquema 
lineal y usando la identidad $e^{ikh}-2+e^{-ikh} = 
-4\sin^2(kh/2)$, definimos:

\begin{equation}
    \mu := \frac{4\,\delta t\sin^2(kh/2)}{h^2} > 0.
\end{equation}

Los tres esquemas dan los siguientes factores de amplificación:

\begin{center}
\renewcommand{\arraystretch}{2.2}
\begin{tabular}{|l|c|c|c|}
\hline
\textbf{Esquema} & \textbf{Factor} $g$ & $|g|$ & 
\textbf{Consecuencia} \\
\hline
Euler explícito & $1+i\mu$ & $\sqrt{1+\mu^2}>1$ & 
$N_\psi^n$ crece, inestable \\
\hline
Euler implícito & $\dfrac{i}{i+\mu}$ & 
$\dfrac{1}{\sqrt{1+\mu^2}}<1$ & 
$N_\psi^n$ decrece, disipativo \\
\hline
Crank-Nicolson & $\dfrac{2i-\mu}{2i+\mu}$ & $1$ & 
$N_\psi^n$ conservada \\
\hline
\end{tabular}
\end{center}

Para Crank-Nicolson, $2i-\mu$ y $2i+\mu$ son complejos 
conjugados — misma parte real $\mu$, partes imaginarias 
opuestas $\pm 2$ — por lo que tienen el mismo módulo:

\begin{equation*}
    |g| = \frac{|2i-\mu|}{|2i+\mu|} = 
    \frac{\sqrt{\mu^2+4}}{\sqrt{\mu^2+4}} = 1.
\end{equation*}

Esto vale para todo $k$, $h$ y $\delta t$, sin ninguna 
restricción. La conservación de la norma discreta sigue 
inmediatamente:

\begin{equation*}
    N_\psi^{n+1} = h\sum_k|c_k|^2|g_k|^{2(n+1)} = 
    h\sum_k|c_k|^2 = N_\psi^0.
\end{equation*}

\subsection{Importancia para los solitones}

Un solitón existe gracias al equilibrio exacto entre 
\textbf{dispersión} ($\partial_{xx}\psi$, que ensancha el 
perfil) y \textbf{no linealidad} ($\psi(1-|\psi|^2-
\alpha|\phi|^2)$, que lo concentra). Este equilibrio 
depende directamente de que $N_\psi$ permanezca constante.

Si el esquema no conserva la norma:

\begin{itemize}
    \item \textbf{Euler explícito} ($N_\psi^n$ crece): el 
    esquema inyecta densidad ficticia paso a paso. La no 
    linealidad domina artificialmente y el solitón parece 
    explotar — \textbf{inestabilidad falsa}.
    \item \textbf{Euler implícito} ($N_\psi^n$ decrece): 
    el esquema extrae densidad paso a paso. La dispersión 
    domina artificialmente y el solitón se aplana y 
    desaparece — \textbf{estabilidad falsa}.
\end{itemize}

Con \textbf{Crank-Nicolson} ($N_\psi^n = N_\psi^0$ 
exactamente), el balance se mantiene correcto en todo 
momento y el resultado de la simulación refleja 
genuinamente la respuesta del solitón a la perturbación. 
Para cuantificar la estabilidad seguimos a Di Menza y 
Gallo y calculamos:

\begin{equation}
    k(t) := \frac{\log(1+\|u(t)\|_{L^\infty})}{t}.
\end{equation}

Si $k(t)\to 0$ cuando $t\to\infty$, no hay crecimiento 
exponencial y el solitón es \textbf{linealmente estable}. 
Si $k(t)\to\lambda>0$, existe un modo inestable con tasa 
de crecimiento $\lambda$. Esta distinción solo es posible 
si el esquema no introduce crecimiento ni decaimiento 
artificiales de la norma, lo que garantiza Crank-Nicolson.
---------------------------------

\textbf{MIS NOTAS}

 Asi una vez discretizado el espacio, la condición de conservación que debe 
respetar el esquema numérico pasa a ser la conservación de la 
\textbf{norma discreta}:

\begin{equation}
    N_\psi^n = h\sum_j|\psi_j^n|^2, \qquad 
    N_\psi^{n+1} = N_\psi^n.
\end{equation}




\end{document}
